%===============================================================================
% LaTeX sjabloon voor de bachelorproef toegepaste informatica aan HOGENT
% Meer info op https://github.com/HoGentTIN/latex-hogent-report
%===============================================================================

\documentclass[dutch,dit,thesis]{hogentreport}

% TODO:
% - If necessary, replace the option `dit`' with your own department!
%   Valid entries are dbo, dbt, dgz, dit, dlo, dog, dsa, soa
% - If you write your thesis in English (remark: only possible after getting
%   explicit approval!), remove the option "dutch," or replace with "english".

\usepackage{lipsum} % For blind text, can be removed after adding actual content
\usepackage{tabularx}
\usepackage{float}    % voeg dit toe aan je preamble
\usepackage{booktabs} % als je dit nog niet hebt
\usepackage[chapter]{minted}
\usemintedstyle{solarized-light}

%% Formatting for minted environments.
\setminted{%
    autogobble,
    frame=lines,
    breaklines,
    linenos,
    tabsize=4
}

\usepackage{listings}
\usepackage{xcolor}

\lstdefinestyle{custompython}{
  language=Python,
  basicstyle=\ttfamily\small\color{black},
  keywordstyle=\color{blue}\bfseries,
  stringstyle=\color{teal},
  commentstyle=\color{gray}\itshape,
  showstringspaces=false,
  breaklines=true,
  frame=single,
  captionpos=b,
  numbers=left,
  numberstyle=\tiny,
}

\definecolor{Rood}{HTML}{E74C3C}
\definecolor{Blauw}{HTML}{3498DB}
\definecolor{Groen}{HTML}{2ECC71}
\definecolor{Oranje}{HTML}{F39C12}
\definecolor{Paars}{HTML}{9B59B6}
\definecolor{Roos}{HTML}{E91E63}
\definecolor{Lime}{HTML}{CDDC39}
\definecolor{Grijs}{HTML}{95A5A6}
\definecolor{Geel}{HTML}{F1C40F}
\definecolor{HogentOranje}{HTML}{EF8767}
\definecolor{HogentLichtOranje}{HTML}{F8CCA2}
\definecolor{HogentBlauw}{HTML}{0000FF}
\definecolor{HogentLichtBlauw}{HTML}{A4C8F0}
%% For source code highlighting, requires pygments to be installed
%% Compile with the -shell-escape flag!
%% \usepackage[chapter]{minted}
%% If you compile with the make_thesis.{bat,sh} script, use the following
%% import instead:
\usepackage[chapter,outputdir=../output]{minted}
\usemintedstyle{friendly}

%% Formatting for minted environments.
\setminted{%
    autogobble,
    frame=lines,
    breaklines,
    linenos,
    tabsize=4
}

%% Ensure the list of listings is in the table of contents
\renewcommand\listoflistingscaption{%
    \IfLanguageName{dutch}{Lijst van codefragmenten}{List of listings}
}
\renewcommand\listingscaption{%
    \IfLanguageName{dutch}{Codefragment}{Listing}
}
\renewcommand*\listoflistings{%
    \cleardoublepage\phantomsection\addcontentsline{toc}{chapter}{\listoflistingscaption}%
    \listof{listing}{\listoflistingscaption}%
}

\usepackage[utf8]{inputenc}
\usepackage{glossaries}
% Other packages not already included can be imported here
% \makeglossaries


%%---------- Document metadata -------------------------------------------------
% TODO: Replace this with your own information
\author{Phil Van Akoleyen }
\supervisor{Dhr. S. Samyen }
\cosupervisor{Dhr. K. Haeck}
\title[Optionele ondertitel]%
    {Machine learning-gebaseerde anomaliedetectie in web-en API-verkeer: PoC binnen de securityomgeving van Evolane.}
\academicyear{\advance\year by -1 \the\year--\advance\year by 1 \the\year}
\examperiod{1}
\degreesought{\IfLanguageName{dutch}{Professionele bachelor in de toegepaste informatica}{Bachelor of applied computer science}}
\partialthesis{false} %% To display 'in partial fulfilment'
%\institution{Internshipcompany BVBA.}


%% Add global exceptions to the hyphenation here
\hyphenation{back-slash}

%% The bibliography (style and settings are  found in hogentthesis.cls)
\addbibresource{bachproef.bib}            %% Bibliography file
\addbibresource{../voorstel/voorstel.bib} %% Bibliography research proposal
\defbibheading{bibempty}{}

%% Prevent empty pages for right-handed chapter starts in twoside mode
\renewcommand{\cleardoublepage}{\clearpage}

\renewcommand{\arraystretch}{1.2}

%% Content starts here.
\begin{document}

%---------- Front matter -------------------------------------------------------

\frontmatter

\hypersetup{pageanchor=false} %% Disable page numbering references
%% Render a Dutch outer title page if the main language is English
\IfLanguageName{english}{%
    %% If necessary, information can be changed here
    \degreesought{Professionele Bachelor toegepaste informatica}%
    \begin{otherlanguage}{dutch}%
       \maketitle%
    \end{otherlanguage}%
}{}

%% Generates title page content
\maketitle
\hypersetup{pageanchor=true}

%%=============================================================================
%% Voorwoord
%%=============================================================================

\chapter*{\IfLanguageName{dutch}{Woord vooraf}{Preface}}%
\label{ch:voorwoord}

%% TODO:
%% Het voorwoord is het enige deel van de bachelorproef waar je vanuit je
%% eigen standpunt (``ik-vorm'') mag schrijven. Je kan hier bv. motiveren
%% waarom jij het onderwerp wil bespreken.
%% Vergeet ook niet te bedanken wie je geholpen/gesteund/... heeft
Deze bachelorproef vormt het laatste onderdeel van mijn opleiding Toegepaste Informatica aan Hogeschool Gent. Het onderzoek bood me de kans om mijn interesse in cybersecurity te combineren met machine learning, een combinatie die me doorheen het schrijven van deze bachelorproef enorm heeft geboeid.
Graag wil ik mijn promotor, dhr. S. Samyen bedanken. Zijn feedback en begeleiding hebben me geholpen om het werk gestructureerd en doelgericht aan te pakken. Daarnaast wil ik mijn dank uiten aan mijn co-promotor, dhr. K. Haeck, voor de kans om mijn bachelorproef binnen Evolane te realiseren. Zijn steun en technische ondersteuning doorheen het volledige proces waren van grote waarde voor dit onderzoek.


\lipsum[1-2]
\input{samenvatting}

%---------- Inhoud, lijst figuren, ... -----------------------------------------

\tableofcontents

% In a list of figures, the complete caption will be included. To prevent this,
% ALWAYS add a short description in the caption!
%
%  \caption[short description]{elaborate description}
%
% If you do, only the short description will be used in the list of figures

\listoffigures

% If you included tables and/or source code listings, uncomment the appropriate
% lines.
\listoftables

\listoflistings

% Als je een lijst van afkortingen of termen wil toevoegen, dan hoort die
% hier thuis. Gebruik bijvoorbeeld de ``glossaries'' package.
% https://www.overleaf.com/learn/latex/Glossaries

%---------- Kern ---------------------------------------------------------------

\mainmatter{}

% De eerste hoofdstukken van een bachelorproef zijn meestal een inleiding op
% het onderwerp, literatuurstudie en verantwoording methodologie.
% Aarzel niet om een meer beschrijvende titel aan deze hoofdstukken te geven of
% om bijvoorbeeld de inleiding en/of stand van zaken over meerdere hoofdstukken
% te verspreiden!

%%=============================================================================
%% Inleiding
%%=============================================================================

\chapter{\IfLanguageName{dutch}{Inleiding}{Introduction}}%
\label{ch:inleiding}

De inleiding moet de lezer net genoeg informatie verschaffen om het onderwerp te begrijpen en in te zien waarom de onderzoeksvraag de moeite waard is om te onderzoeken. In de inleiding ga je literatuurverwijzingen beperken, zodat de tekst vlot leesbaar blijft. Je kan de inleiding verder onderverdelen in secties als dit de tekst verduidelijkt. Zaken die aan bod kunnen komen in de inleiding~\autocite{Pollefliet2011}:

\begin{itemize}
  \item context, achtergrond
  \item afbakenen van het onderwerp
  \item verantwoording van het onderwerp, methodologie
  \item probleemstelling
  \item onderzoeksdoelstelling
  \item onderzoeksvraag
  \item \ldots
\end{itemize}

\section{\IfLanguageName{dutch}{Probleemstelling}{Problem Statement}}%
\label{sec:probleemstelling}

Uit je probleemstelling moet duidelijk zijn dat je onderzoek een meerwaarde heeft voor een concrete doelgroep. De doelgroep moet goed gedefinieerd en afgelijnd zijn. Doelgroepen als ``bedrijven,'' ``KMO's'', systeembeheerders, enz.~zijn nog te vaag. Als je een lijstje kan maken van de personen/organisaties die een meerwaarde zullen vinden in deze bachelorproef (dit is eigenlijk je steekproefkader), dan is dat een indicatie dat de doelgroep goed gedefinieerd is. Dit kan een enkel bedrijf zijn of zelfs één persoon (je co-promotor/opdrachtgever).

\section{\IfLanguageName{dutch}{Onderzoeksvraag}{Research question}}%
\label{sec:onderzoeksvraag}

Hoe kan een machine learning-gedreven anomaliedetectiesysteem op basis van Akamai access logs worden ontworpen en geïmplementeerd om binnen de securityomgeving van Evolane afwijkend web- en API-verkeer te detecteren en het aantal false positives te reduceren? 

\section{\IfLanguageName{dutch}{Onderzoeksdoelstelling}{Research objective}}%
\label{sec:onderzoeksdoelstelling}

Het doel van deze bachelorproef is het ontwikkelen van een werkend proof-of-concept van een machine learning-gedreven anomaliedetectiesysteem dat afwijkend gedrag in web- en API-verkeer detecteert op basis van Akamai access logs, binnen de securityomgeving van Evolane. 

Het prototype is geslaagd wanneer het aantoonbaar is dat het anomalieën kan identificeren op basis van gedragsafwijking en er een meetbare vermindering is in het aantal false positives ten opzichte van de huidige benadering. De evaluatie hiervan gebeurt aan de hand van volgende meetcriteria: F1-score, precision, recall en false positive ratio. Naast de technische evaluatie wordt ook de operationele bruikbaarheid beoordeeld, specifieker of het integreerbaar is met de bestaande securityomgeving van Evolane. 
Het eindresultaat bevat naast het proof-of-concept ook een onderbouwd advies voor Evolane over de toepasbaarheid, voordelen, beperkingen en de mogelijke vervolgstappen voor een ML-gedreven anomaliedetectie-systeem binnen hun omgeving.  

\section{\IfLanguageName{dutch}{Opzet van deze bachelorproef}{Structure of this bachelor thesis}}%
\label{sec:opzet-bachelorproef}

% Het is gebruikelijk aan het einde van de inleiding een overzicht te
% geven van de opbouw van de rest van de tekst. Deze sectie bevat al een aanzet
% die je kan aanvullen/aanpassen in functie van je eigen tekst.

De rest van deze bachelorproef is als volgt opgebouwd:

In Hoofdstuk~\ref{ch:stand-van-zaken} wordt een overzicht gegeven van de stand van zaken binnen het onderzoeksdomein, op basis van een literatuurstudie.

In Hoofdstuk~\ref{ch:methodologie} wordt de methodologie toegelicht en worden de gebruikte onderzoekstechnieken besproken om een antwoord te kunnen formuleren op de onderzoeksvragen.

% TODO: Vul hier aan voor je eigen hoofstukken, één of twee zinnen per hoofdstuk

In Hoofdstuk~\ref{ch:conclusie}, tenslotte, wordt de conclusie gegeven en een antwoord geformuleerd op de onderzoeksvragen. Daarbij wordt ook een aanzet gegeven voor toekomstig onderzoek binnen dit domein.
\chapter{\IfLanguageName{dutch}{Stand van zaken}{State of the art}}%
\label{ch:stand-van-zaken}

% Eigen stand van zaken: 
\section{OWASP en hedendaagse web- en API-dreigingen}
\label{sec:api-security-context}
Om de rol van anomaliedetectie binnen webbeveiliging te kaderen, is het belangrijk eerst de huidige dreigingen te schetsen en de beperkingen van bestaande detectiemechanismen te bespreken. 

Binnen web- en API-beveiliging speelt OWASP, het Open Web Application Security Project, een centrale rol voor het identificeren van de tien meest kritieke kwetsbaarheden. OWASP publiceert periodiek top tien lijsten voor zowel webapplicaties als voor API's\autocite{OWASP2026,OWASP2023}. Deze lijsten bieden een overzicht van de meest voorkomende en impactvolle kwetsbaarheden binnen het domein van webbeveiliging.%OWASP WEB EN API BRON. 

Voor webbapplicaties omvat de OWASP TOP 10 onder meer kwetsbaarheden zoals injection-aanvallen, broken authentication en security misconfigurations. Deze kwetsbaarheden zijn voornamelijk gericht tot de interactie tussen gebruiker en de webinterface\autocite{OWASP2026}%OWASP WEB BRON. 

Voor API's verschuift de focus naar andere soorten risico's zoals Broken Object Level Authorization (BOLA), Broken Function Level Authorization en Unrestricted Resource Consumption. API kwetsbaarheden zijn vaker gerelateerd aan onvoldoende controle op de toegang tot deze APIs en het misbruiken van businesslogica\autocite{OWASP2023}%API BRON. 

Het onderscheid dat gemaakt wordt tussen web- en API-kwetsbaarheden benadrukt dat er nood is aan beveiliging op meerdere niveaus. Waar dat webaanvallen eerder gericht zijn op de gebruikersinterface, richten API-aanvallen zich tot de directe communicatie tussen systemen. Dit onderscheid is de basis voor een diepere analyse van hoe aanvallen op beide niveaus evolueren. 

\subsection{De groeiende complexiteit van Web-aanvallen}
\label{subsec:api-aanvallen-complexiteit}
De complexiteit van aanvallen gericht op webapplicaties is in de afgelopen jaren sterk toegenomen. Volgens rapporten van Akamai is er sprake een significante stijging in webaanvallen, hier spreken ze onder andere over een toename van 33 procent tussen 2023 en 2024\autocite{Akamai2025}%Akamai bron.

Binnen de OWASP TOP 10 komen verschillende aanvalsvormen naar voren die deze evolutie illustreren. Een belangrijk voorbeeld is Broken Access Control waar een aanvaller buiten zijn toegestane rechten handelingen uitvoert. Daarnaast blijven injection attacks een voorkomend risico, dit is wanneer de aanvaller code zoals SQL of OS commando's via de webapplicatie uit laat voeren op de server waar de applicatie gehost word om zo data te verkrijgen. Ook authentication failures vormen een ketsbaarheid, dit is wanneer de beveiliging van de applicatie er onvoldoende in slaagt om de user van de aanvaller te detecteren als ongeldig of incorrect dit kan dan leiden tot brute force en geautomatiseerde aanvallen\autocite{OWASP2026}%bron owasp WEB.

Een kritieke evolutie binnen deze aanvalsvormen is dat kwaadaardig verkeer steeds moeilijker te onderscheiden is van normaal geldig verkeer. Vandaag de dag maken aanvallers steeds vaker gebruik van geautomatiseerde bots. Deze bots zijn speciaal ontworpen om menselijk gedrag zo nauwkeurig mogelijk te simuleren door het implementeren van variatie in de frequentie van verzonden requests en hun navigatiepatronen. Deze geautomatiseerde aanvallen maken gebruik van grote aantallen requests die geldig lijken en zich aan kunnen passen naar de beveiliging van een webapplicatie. Dit zorgt ervoor dat het verschil tussen normaal en kwaadaardig verkeer moeilijker waarneembaar wordt\autocite{F52026}.%F5 bron 

Naast de stijgende complexiteit van webaanvallen kent ook het API-landschap een gelijkaardige evolutie. 
\subsection{De groeiende complexiteit van API-aanvallen}
\label{subsec:api-aanvallen-complexiteit}

De complexiteit en frequentie van aanvallen op APIs is eveneens sterk toegenomen. APIs vormen een essentieel onderdeel voor de moderne webapplicaties, maar ze vergroten het aanvalsoppervlak aanzienlijk\autocite{Akamai2025}. %Akamai bron

In tegenstelling tot klassieke webaanvallen richten API-aanvallen zich op de communicatie tussen softwarecomponenten. Volgens OWASP ontstaan API-kwetsbaarheden voornamelijk door foutieve of beperkte authenticatie en autorisatie. Kwetsbaarheden zoals Broken Object Level Authorization en Broken Function Level Authorization tonen aan dat aanvallers zich richten tot endpoints die technisch correct functioneren, maar onvoldoende toegangscontrole bevatten.

Daarnaast worden API’s vaak misbruikt via geautomatiseerde requests, bijvoorbeeld bij Unrestricted Resource Consumption en business logic abuse. Dit type aanvallen maakt gebruik van legitieme interacties, waardoor detectie via traditionele signature-based systemen moeilijk is\autocite{OWASP2023}.%bron OWASP 

Zowel bij webaanvallen als bij API-aanvallen stelt zich dezelfde vraag: waarom slagen bestaande beveiligingsmechanismen er onvoldoende in om deze dreigingen tijdig te detecteren. 

\subsection{Beperkingen van traditionele detectiemechanismen}
\label{subsec:traditioneel-beperkingen}
% Leg uit waarom rule-based systemen tekorschieten bij moderne dreigingen.

Traditionele beveiligingsmechanismen zoals Web Application Firewall (WAF), signature based Intrusion Detection Systems (IDS) en andere regelgebaseerde detectiesystemen steunen voornamelijk op bekende aanvalssignaturen en patronen die op voorhand gedefineerd zijn. Zoals \textcite{Standards2007} beschrijft in hun Intrusion Detection and Prevention systems richtlijnen zijn deze systemen zeer effectief in het detecteren van gedocumenteerde dreigingen, maar eerder beperkt in het detecteren van nieuwe aanvallen waarvoor nog geen signature aanval bestaat. 

Zoals eerder beschreven maken moderne aanvallers steeds vaker gebruik van geautomatiseerde technieken die gespecialiseerd zijn in het nabootsen van legitiem verkeer en dat zich aanpast aan de beveiliging van een applicatie. Dergelijke aanvallen bevatten vaak geen gekende signatuur of vertonen geen afwijkend patroon dat op voorhand gedifinieerd kan worden \autocite{F52026}.

Recente literatuur bevestigt deze claims en benadrukt signature-based systemen te beperkt zijn in het detecteren van zero-day aanvallen binnen het snel evoluerende cybersecuritylandschap. Daarom wordt machine learning steeds vaker voorgesteld als een aanvulling op deze bestaande systemen om eerder ongekende aanvallen te detecteren, false positives te verminderen en beveiliging te verbeteren\autocite{Kumar_2025}. % komt van deze bron Impact of Machine Learning on Intrusion Detection Systems for the Protection of Critical Infrastructure


\section{Anomaliedetectie binnen Security Operations}
\label{sec:anomaliedetectie-soc}
%Zou hier graag nog een subsectie aan toe willen voegen

De beperkingen van traditionele detectiemechanismen, zoals beschreven in \ref{subsec:traditioneel-beperkingen}, creëren een behoefte voor een aanvullend detectiesysteem binnen het Security Operations Team.  Binnen de context van web- en API-beveiliging speelt het SOC een centrale rol in het bewaken van inkomend verkeer en het detecteren van verdachte events.

\subsection{Akamai APP en API-protector}
\label{subsec:App-API}
Akamai App \& API Protector is een geïntegreerde beveiligingsoplossing die webbapplicaties en API's beschermt tegen een breed aantal dreigingen, waaronder Dos-aanvallen, botactiviteit en gekende kwetsbaarheden uit onder andere de OWASP top tien. Dit doet het door het combineren van meerdere modules waaronder een WAF, botmitigatie en API-bescherming in één platform.%https://www.akamai.com/products/app-and-api-protector

De bescherming van webapplicaties gebeurt via een regelgebaseerde evaluatie van inkomende requests. Indien een request voldoet aan vooraf gedefinieerde criteria, past het platform een set beveiligingscontroles toe. Indien een aanvalsvector wordt gedetecteerd, wordt de toegang geblokkeerd of wordt er een alert gegenereerd afhankelijk van de geconfigureerde instellingen\autocite{Technologies2026}.

Voor botdetectie maakt het platform gebruik van Bot-manager-module, die naast regelgebaseerde signaturen ook machine learning inzet voor het detecteren van bots. Deze module richt zich op een breed spectrum aan detectiemethoden: browser fingerprinting, request-anomalieën, reputatiescoring en gedragsanalyse. De Bot manager is specifiek ontworpen voor het indentificeren van bots binnen het verkeer tussen de menselijke gebruikers.%https://www.akamai.com/products/bot-manager 

Voor API-bescherming biedt het platform zichtbaarheid over alle API-endpoints binnen een omgeving, inclusief niet-gedocumenteerde of zogenaamde shadow  API's. Het platform analyseert API-verkeer op basis van vooraf bepaalde beleidsregels en detecteert afwijkingen zoals buitensporige requests of ongeautoriseerde toegangspogingen\autocite{Akamai2025}.
\subsection{De toegevoegde waarde van ML in SOC-context}
\label{subsec:toegevoegde-waarde-ML}
% Motiveer waarom ML een waardevolle aanvulling is op klassieke detectiesysteem bron: https://dl.acm.org/doi/10.1145/3723158?
Een Security Operations Center (SOC) is verantwoordelijk voor het continu monitoren en analyseren van web- en API-verkeer op zoek naar verdachte activiteit en beveiligingsincidenten. Binnen een webbeveiligingscontext maken SOC-analisten gebruik van tools zoals WAF-systemen, IDS-systemen en SIEM-platformen (Security Information and Event Management) om inkomende requests te evalueren en potentiële aanvallen tijdig te detecteren en te beantwoorden.

In moderne Security Operations Teams leidt een grote hoeveelheid security alerts tot een fenomeen wat door de literatuur beschreven wordt als alert fatigue. Dit is wanneer analisten worden overweldigd door een constante stroom van alerts en false positives. Met als gevolg een verhoogde operationele belasting, wat het risico biedt dat kritieke incidenten worden gemist. 

De toegevoegde waarde van machine learning binnen een SOC-team heeft voordelen op verschillende niveaus: na het inwerken van het model een reductie van false positives, verbeterde detectie van abnormaal gedrag en het schakelen van reactief handelen naar proactief detecteren\textcite{Miranda2025}. %https://www.crowdstrike.com/en-us/cybersecurity-101/next-gen-siem/ai-siem/


\section{Machine learning-methoden voor anomaliedetectie}
\label{sec:ML-methoden-anomalie}

Zoals beschreven in sectie \ref{subsec:traditioneel-beperkingen} hebben de beperkingen van traditionele rule-based detectiesystemen ervoor gezorgd dat de interesse in machine learning gedreven benaderingen voor anomaliedetectie zijn gestegen. Deze sectie biedt een overzicht van de belangrijkste machine learning gedreven benaderingen voor anomaliedetectie, met specifieke aandacht op hun toepasbaarheid. 

\subsection{supervised, unsupervised en semi-supervised}
\label{subsec:toegevoegde-waarde-ML}
% overzicht geven van de 3 hoofdbenaderingen met voor- en nadelen. 

Machine learning benaderingen voor anomaliedetectie worden onderverdeeld in 3 categorieën op basis van de nood aan gelabelde trainingsdata. \textcite{Saabith2023} maakt een onderscheid in drie hoofdcategorieën: supervised learning, unsupervised learning en semi-supervised learning. Deze benaderingen bevatten elk sterke punten en zwakke punten die hun toepasbaarheid bepalen binnen security context.


\subsubsection{Supervised learning}

Supervised learning-algoritmen vereisen een dataset waarbij elk datapunt uitdrukkelijk gelabeld is als normaal of abnormaal gedrag, vervolgens leert het model tijdens de trainingsfase de relatie tussen de invoerkenmerken (features in de dataset) en de daarbij horende labels hierdoor is het instaat om nieuwe datapunten te classificeren~\autocite{Saabith2023}. 
De algoritmen binnen deze categorie zijn: bayesian networks, k-nearest neighbors, random forest, supervised neural networks en suport vector machines. Over het algemeen ligt de detectierate hoger bij supervised modellen dan bij unsupervised modellen, maar dit gaat dan ook gepaard met de moeilijkheidsgraad in het trainen van het model waarbij dat bij supervised hoger ligt\autocite{Kumari2024}. 
De voornaamste beperking van supervised learning-algoritmenn is de afhankelijkheid aan een gelabelde dataset. In werkelijkheid zijn gelabelde datasets weinig te vinden. Ten eerste vereist het labelen van data security-analisten met expertise in het domein dat duizenden events dienen te categoriseren, een proces dat lang duurt~\autocite{Kumar_2025}. 
Vervolgens kan er niet worden gegarandeerd dat de trainingsdata alle relevante aanvalstypes bevat, zo is het mogelijk dat niet bekende aanvalstypes, zero-day exploits en andere varianten van aanvallen mogelijks niet aanwezig zijn in de trainingsdata van het model. Dit kan ervoor zorgen dat het model nieuwe subtiele anomalieën mist~\autocite{Kumari2024}.

\subsubsection{Unsupervised learning}

Unsupervised learning-algoritmen zijn in tegenstelling tot supervised benaderingen niet gebonden aan een gelabelde dataset waarop het model getraind moet worden. Deze modellen leren om patronen en structuren te herkennen in de data op basis van zowel gelijkheden als verschillen tussen datapunten. Punten die van het opgebouwde patroon afwijken of niet binnen een bestaande vooraf gedefinieerde  cluster past zullen worden gezien als een anomalie\textcite{Saabith2023}. 
Algoritmen binnen deze categorie zijn: clustering methoden (K-means, DBSCAN), isolation forest, local outlier factor, Principal Component Analysis en Autoencoders. een van de voordelen van unsupervised machine learning modellen is dat ze uitblinken in het detecteren van ongekende anomalieën doordat ze zich niet baseren op historische trainingsdata\autocite{Kumari2024}.

Een andere nuttige eigenschap van unsupervised learning is dat deze modellen continue leren wat normaal is op basis van de binnenkomende data, dit zorgt ervoor dat ze schaalbaar zijn in omgevingen waar veel data binnenkomt en ze zich kunnen aanpassen naarmate de data veranderd\autocite{Traceable2024}.

De voornaamste beperking van unsupervised learning is dat het kan resulteren in een hogere false positive rate ten opzichte van een supervised benadering. Doordat het model geen expliciete instructies ontvangen heeft zal het alle statistische datapunten die uitschieten beschouwen als verdacht, dit kan dus ook zeldzaam legitiem verkeer bevatten\autocite{Saabith2023}. 


\subsubsection{Semi supervised learning}

Semi-supervised learning is een combinatie benadering van beide voorgaande methoden, het model wordt dus getraind op zowel niet gelabelde data als gelabelde data. De unlabeld data dienen om de accuraatheid van het model te verhogen op niet gekende anomaliteiten en de labeld data geeft het model zicht op welke afwijkingen er voornamelijk zijn\autocite{Saabith2023}. 
Algoritmen binnen deze categorie zijn one-class SVM en autoencoder met training data. Autoencoder kan zowel usnupervised als semi-supervised dit is afhankelijk van het wel of niet implementeren van labeld data\autocite{Bajallan2020}. %komt uit deze bron https://kth.diva-portal.org/smash/get/diva2:1465807/FULLTEXT01.pdf
Doordat semi-supervised beperkte labeld data gebruikt kan het accuratere beslissingen maken op wat wel of niet een afwijking is op het normale gedrag. Hierdoor kan het ten opzichte van een unsupervised benadering de false positive rate reduceren\autocite{EmergentMind_2026}.  % komt uit https://www.emergentmind.com/topics/semi-supervised-anomaly-detection bron 
De voornaamste restrictie van een semi-supervised benadering is dat er hoewel in mindere maten nog steeds nood is aan gelabelde data voor het trainen van het model \autocite{Saabith2023}.

\subsection{Veelgebruikte algoritmen per categorie}%
\label{subsec:ml-algoritmen-overzicht}

\subsubsection{Supervised ML modellen}

\paragraph{Random Forest}

Random Forest is een ensemble-gebaseerde supervised machine learning algoritme dat gebruikt kan worden voor anomaliedetectie. 
Het model bestaat uit een verzameling van decision trees, waarbij elke boom getrained wordt op een willekeurige subset uit de trainingsdata en een willekeurige set van features gebruikt\autocite{Breiman2001}.

Random Forest kan herleid worden tot een verzameling boomgebaseerde classifiers, hierbij is elke boom opgebouwd door een onafhankelijke en identieke vector die willekeurig gekozen is door het model. 
Voor een gegeven inputvector zal het model een voorspelling geven, waarna de classificatie bepaald wordt door majority voting.
In de context van anomaliedetectie zal de boom een classificatie geven zoals normaal of anomaal, en de uiteindelijke beslissing is gebaseerd op een meerderheid in stemmen\autocite{Primartha2017}.

Binnen Random Forest zijn er enkele essentiële parameters die direct invloed hebben op de werking en stabiliteit van het model. 
De \texttt{n\_estimators} parameter staat in voor het aantal decision trees binnen het ensemble; een hoger aantal bomen verhoogt de stabiliteit maar heeft ook rechtstreeks invloed op de rekencomplexiteit. 
De \texttt{max\_features} parameter bepaalt hoeveel features willekeurig gekozen worden per splitsing in de decision tree\autocite{Breiman2001}.

Een fundamentele beperking van Random Forest is de nood aan gelabelde trainings data
waarbij elk datapunt in de dataset gelabeld moet zijn als normaal of anomaal. Dit maakt Random Forest minder geschikt voor omgevingen waarin de aard van het verkeer niet op voorhand gekend is en de nadruk ligt op het ontdekken van voorheen onbekende anomalieën~\autocite{Breiman2001}

\textbf{Voordelen:}
\begin{itemize}
    \item Robuust tegen overfitting, doordat het model niet sterk afstemt op
    zijn trainingsdata en goed generaliseert naar nieuwe, ongeziene data.
    \item Expliciete featureselectie is doorgaans niet noodzakelijk, doordat
    het model efficiënt omgaat met een groot aantal variabelen.
\end{itemize}

\textbf{Nadelen:}
\begin{itemize}
    \item Lage interpreteerbaarheid, doordat de uiteindelijke voorspelling
    gemaakt wordt door een combinatie van decision trees, is het moeilijk om
    de grondslag van een classificatie te achterhalen.
    \item Prestaties kunnen dalen bij sterk gecorreleerde features, wat de
    diversiteit tussen de decision trees vermindert.
\end{itemize}

\medskip

\paragraph{Support Vector Machine}

Support Vector Machine is een supervised classificatiealgoritme dat een scheidingslijn genaamd de hyperplane trekt tussen twee klassen. 
Bij binaire classificatie zoals het onderscheid maken tussen normaal en anomalie, probeert SVM de grens te vinden die de twee klassen scheidt met een zo groot mogelijke marge. 
De marge is de afstand tussen de beslissingsgrens en de datapunten van beide klassen die zich het dichtst bij bevinden, deze datapunten worden ook welde support vectors genoemd. 
Door de marge zo groot mogelijk te maken kan SVM beter generaliseren naar nieuwe data en is het minder gevoelig aan kleine variaties in de set\autocite{Cortes1995}.

In tegenstelling tot de praktijk is hier data zelden linair scheidbaar: de grens tussen normaal of anomaal kan zelden met een rechte lijn worden getrokken. 
Om deze uitdaging aan te pakken maakt SVM gebruik van kernelfuncties. Een kernelfunctie transformeert data richting een hogere dimensionale ruimte waarin de data dan weer wel linair scheidbaar is, dit wordt ook wel de kernel trick genoemd.
Het gebruik van een kernelfunctie heeft ook een rechtstreeks positieve invloed op de berekening van de transformatie. 
SVM gebruikt voornamelijk kernelfuncties zoals de lineaire kernel, de polynomiale kernel en de Radial Basis Function kernel. De keuze van kernel is sterk afhankelijk met de structuur van de data. %BRON: 

\textbf{Voordelen:}
\begin{itemize}
    \item Sterk in hoog-dimensionale datasets, aangezien het model goed overweg
    kan met een groot aantal features binnen de data.
    \item Capabel model voor het aanduiden van grenzen tussen klassen in de data.
\end{itemize}

\textbf{Nadelen:}
\begin{itemize}
    \item Lage interpreteerbaarheid, het is moeilijk om te achterhalen wat de
    grondslag was van een bepaalde classificatie.
    \item Gevoelig aan parameter settings en trainingstijd kan exponentieel
    stijgen bij grote datasets.
\end{itemize}

\medskip

\subsubsection{Unsupervised ML modellen}

\paragraph{Isolation Forest}

Isolation Forest is een unsupervised anomaliedetectiemethode dat vertrekt vanuit een andere invalshoek ten opzichte van andere ML modellen: in plaats van een profiel op te bouwen van hoe normaal gedrag eruitziet, focust Isolation Forest zich op het isoleren van afwijkende datapunten. 
Dit is mogelijk doordat het model ervan uitgaat dat anomalieën doorgaans zeldzaam zijn en afwijkende waarden bezitten\autocite{Liu2008}.

Het algoritme bouwt een vooraf gedefinieerd aantal binaire decision trees. 
In elke boom wordt iteratief een willekeurige feature en een willekeurige splitswaarde geselecteerd om datapunten te scheiden. 
De nodige padlengte om een datapunt volledig te isoleren wordt gemeten en genormaliseerd tot een anomaliescore tussen 0 en -1.
Datapunten met een korte gemiddelde padlengte over alle bomen heen krijgen een hoge anomaliescore\autocite{Lesouple2021}.

Isolation Forest maakt gebruik van subsampeling voor het trainen van de decision trees binnen het model, dit wil zeggen dat niet elke boom getrained wordt op de volledige dataset maar op een willekeurige subset van de data.
Dit verbetert de detectiekwaliteit van het model doordat het 2 fenomenen tegengaat. 
Enerzijds voorkomt het masking, dit is wanneer datapunten die dicht bij elkaar liggen in een grote dataset onterecht door het model zouden worden beschouwd als een kleine groep legitieme data. 
Anderzijds voorkomt het swamping, dit is het tegenovergestelde van masking, hierbij worden datapunten die zicht dicht tegen de rand bevinden verward als anomaal. 
Een ander voordeel van subsampeling is dat het model efficiënt schaalt naar grote datasets doordat de trainingstijd per boom beperkt blijft. 
Bij de training van het model dient een contamination-parameter te worden gedefineerd. Deze parameter is verantwoordelijk voor het bepalen welk percentage van de data het model zal classificeren als afwijkend.
De keuze van deze parameter is afhankelijk van de context van het model en de data en heeft rechtstreekt invloed op de balans tussen het detecteren van normaal of anomaal verkeer\autocite{Liu2008}. 

\textbf{Voordelen:}
\begin{itemize}
    \item Zeer efficiënt en schaalbaar, doordat het algoritme een lineaire
    tijdcomplexiteit heeft en weinig geheugen vereist, kan het goed toegepast
    worden op grote datasets en data met veel dimensies.
    \item Effectief in het detecteren van anomalieën, omdat het model anomalieën
    direct probeert te isoleren in plaats van normale data te modelleren,
    waardoor afwijkende punten sneller herkend worden.
\end{itemize}

\textbf{Nadelen:}
\begin{itemize}
    \item Moeilijkheden bij overlappende of dichte clusters, wanneer anomalieën
    dicht bij normale data liggen of in grote groepen voorkomen, worden ze minder
    goed geïsoleerd en dus moeilijker gedetecteerd.
    \item Gevoelig voor data-eigenschappen zoals dimensionaliteit en sample size,
    waardoor prestaties kunnen afnemen bij veel irrelevante features of een
    ongunstige keuze van de subsample grootte.
\end{itemize}

\medskip

\paragraph{Local Outlier Factor}

Local Outlier Factor is een unsupervised anomaliedetectiemethode die afwijkende datapunten detecteert op basis van dichtheidsverschillen. LOF evalueert een datapunt relatief ten opzichte van zijn k-dichtste buren.

Het algoritme start met het berekenen van de \texttt{k-distance}, dit is de afstand tussen het datapunt en de k-de dichtste buur. 
Vervolgens wordt de reachability distance berekend, een metric die instabiliteit in de meting voorkomt zodat twee punten die dicht bij elkaar liggen niet automatisch als normaal worden beschouwd. 
De uiteindelijke LOF-score is de verhouding tussen de gemiddelde lokale dichtheid van de buren van het datapunt en de lokale dichtheid van het punt zelf. 
Wanneer de LOF-waarde zakt onder 1 wijst dit op een lagere dichtheid ten opzichte van andere punten en dus een mogelijke anomalie\autocite{Mavuduru2026}.

De voornaamste parameter van LOF is \texttt{k}, het aantal buren dat het model evalueert. 
Een kleinere k-waarde maakt het model gevoeliger voor lokale variaties, een grotere k-waarde maakt het model beter bestendig tegen false positives maar ook minder gevoelig voor lokale anomalieën.

Het voornaamste verschil met isolation forest is dat LOF anomalieën detecteert op basis van relatieve dichtheidsverschillen, dit maakt het model effectiever in het detecteren van lokale anomalieën die zich dichter bevinden bij normale clusters. 
De berekining van de k-nearest neighbors voor elk datapunt schaalt kwadratisch waardoor het minder geschikt is voor grote datasets in een real-time omgeving~\autocite{Breunig2000}. 

\textbf{Voordelen:}
\begin{itemize}
    \item Effectief in het detecteren van lokale anomalieën, doordat het model
    rekening houdt met de dichtheid van naburige datapunten en afwijkingen
    identificeert ten opzichte van de directe omgeving.
    \item Geen gelabelde data vereist, aangezien LOF een unsupervised algoritme
    is dat anomalieën detecteert op basis van structurele eigenschappen van de
    dataset.
\end{itemize}

\textbf{Nadelen:}
\begin{itemize}
    \item Gevoelig voor parameterkeuze, in het bijzonder het aantal buren
    (\texttt{k}), wat een grote invloed heeft op de prestaties en kan leiden
    tot foutieve classificaties.
    \item Kans op false positives aan de randen van clusters, doordat punten
    met een lagere dichtheid aan de grens van normale clusters ten onrechte als
    anomalieën kunnen worden beschouwd.
\end{itemize}

\medskip

\subsubsection{Semi-supervised ML modellen}

\paragraph{One-Class SVM}

One-Class SVM is een variant op de klassieke SVM waarbij het model uitsluitend
getraind wordt op data die als normaal wordt beschouwd. Het doel is om een
beslissingsgrens te leren die de normale data zo goed mogelijk omvat en
afwijkende punten uitsluit.

In tegenstelling tot de klassieke SVM, die twee klassen van elkaar scheidt, zal OC-SVM proberen om een hyperplane te vinden die de data scheidt van de oorsprong.
Doordat datapunten niet altijd direct scheidbaar zijn, maakt SVM gebruik van een kernelfunctie die data projecteert naar een hogere dimensionale ruimte waardoor scheiding mogelijk wordt. 
Het model gebruikt een parameter die de balans bepaalt tussen de toegestane outliers en de complexiteit van het model: enerzijds bepaalt deze een bovengrens voor het aantal toegestane anomalieën, anderzijds een ondergrens voor de toegestane hoeveelheid support vectors\autocite{Hejazi2013}.

De belangrijkste hyperparameter bij one-class SVM is de $\nu$-parameter (\textit{nu}), een waarde tussen 0 en 1 met 2 belangrijke functies. 
Het dient als een bovengrens op het onderdeel aan trainingsfouten, dit zijn de outliers binnen de trainingsdata, en eveneens als een ondergrens op het onderdeel van de support vectors.
Een lagere $\nu$-waarde kan resulteren in een striktere beslissingsgrens die minder outliers toestaat, terwijl een hogere waarde eerder het effect kan hebben dat uitschieters worden geaccepteerd~\autocite{Scholkopf2001}. 

\textbf{Voordelen:}
\begin{itemize}
    \item Geschikt voor anomaliedetectie met één enkele klasse, doordat het model
    een grens leert rond normaal gedrag zonder nood aan expliciet gelabelde
    anomalieën.
    \item Goede prestaties bij imbalanced datasets, aangezien het model ontworpen
    is om zeldzame afwijkingen te detecteren en een lage generalization error kan
    bereiken bij correcte parameterinstelling.
\end{itemize}

\textbf{Nadelen:}
\begin{itemize}
    \item Sterke afhankelijkheid van hyperparameters en kernelkeuze, waarbij een
    verkeerde configuratie kan leiden tot suboptimale prestaties.
    \item Kans op underfitting bij een ongeschikte parameterinstelling, wat
    resulteert in een hoge foutmarge op zowel trainings- als testdata.
\end{itemize}

\medskip

\paragraph{Autoencoder}


Een autoencoder is een neuraal netwerk bestaande uit 2 componenten: een encoder en een decoder. Deze twee componenten zijn verbonden door een tussenlaag genaamd latent space. 
De encoder comprimeert de input in het model naar een compacte representatie in deze latent space, dit zorgt ervoor dat enkel de meest essentiële kenmerken van de data worden behouden. 
Vervolgens probeert de decoder vanuit deze compacte representatie de oorspronkelijke output zo nauwkeurig mogelijk te reconstrueren. Het model wordt getrained de reconstructiefout, het verschil tussen originele input en gereconstrueerde output, te minimaliseren\autocite{Chikezie2025}.

Autoencoders worden in de context van anomaliedetectie voornamelijk op een semi-supervised manier geïmplementeerd. Doordat het model getraind wordt op normale input, is het instaat om gelijkaardige input te reconstrueren.
Echter wanneer input het model passeert dat eerder afwijkend is van het geleerde patroon, zal de reconstructiefout aanzienlijk hoger liggen doordat in de latent space niet de juiste representatie gevormd kon worden. Deze reconstructiefout is de basis als indicator voor een mogelijk anomalie\autocite{Chikezie2025}. 


Een datapunt wordt als anomalie beschouwd wanneer de reconstructiefout de vooraf bepaalde drempelwaarde overschrijdt. Een te lage drempelwaarde kan leiden tot een overvloed aan false positives, terwijl een te hoge waarde mogelijks anomalieën zal missen\autocite{Angiulli2023}.


\textbf{Voordelen:}
\begin{itemize}
    \item Breed toepasbaar en populair in diverse domeinen zoals netwerkdetectie,
    industriële inspectie en beeldanalyse, wat aantoont dat autoencoders flexibel
    inzetbaar zijn voor anomaliedetectie.
    \item In staat om complexe, niet-lineaire patronen in data te modelleren,
    waardoor ze geschikt zijn voor hoog-dimensionale en complexe datasets.
\end{itemize}

\textbf{Nadelen:}
\begin{itemize}
    \item Mogelijkheid tot \textit{out-of-bounds reconstruction}, waarbij
    datapunten ver buiten de trainingsdata toch een lage reconstructiefout
    krijgen, wat leidt tot false negatives.
    \item Gevoeligheid voor generalisatie-effecten zoals interpolatie en
    extrapolatie, waardoor anomalieën kunnen overlappen met normale data en
    onopgemerkt blijven\autocite{Bouman2025}.
    \item De keuze van de netwerk-architectuur, waaronder het aantal lagen en de dimensie van de bottleneck-laag, vereist experimentatie en afstemming. Een te grote bottleneck kan ertoe leiden dat het model ook afwijkende data goed reconstrueert, waardoor anomalieën worden gemist \autocite{Chikezie2025}.
\end{itemize}
\clearpage
\begin{table}[H]
    \centering
    \caption{Vergelijking van machine learning modellen voor anomaliedetectie}
    \label{tab:ml-modellen-vergelijking}
    \begin{tabular}{p{3cm} p{2.8cm} p{4.5cm} p{4.5cm}}
    \toprule
    \textbf{Model} & \textbf{Categorie} & \textbf{Voordelen} & \textbf{Nadelen} \\
    \midrule
    Random Forest 
      & Supervised 
      & Robuust tegen overfitting; geen expliciete featureselectie nodig 
      & Lage interpreteerbaarheid; prestaties dalen bij sterk gecorreleerde features \\
    \addlinespace
    Support Vector Machine 
      & Supervised 
      & Sterk in hoog-dimensionale data; capabel voor klassengrenzen 
      & Lage interpreteerbaarheid; trainingstijd stijgt exponentieel bij grote datasets \\
    \addlinespace
    Isolation Forest 
      & Unsupervised 
      & Lineaire tijdcomplexiteit; effectief door directe isolatie van anomalieën 
      & Moeilijkheden bij overlappende clusters; gevoelig voor dimensionaliteit \\
    \addlinespace
    Local Outlier Factor 
      & Unsupervised 
      & Detecteert lokale anomalieën; geen gelabelde data vereist 
      & Gevoelig voor parameterkeuze (k); kans op false positives aan clusterranden \\
    \addlinespace
    One-Class SVM 
      & Semi-supervised 
      & Geschikt voor één klasse; goede prestaties bij imbalanced datasets 
      & Sterke afhankelijkheid van hyperparameters; kans op underfitting \\
    \addlinespace
    Autoencoder 
      & Semi-supervised 
      & Breed toepasbaar; modelleert complexe niet-lineaire patronen 
      & Risico op out-of-bounds reconstructie; gevoelig voor generalisatie-effecten \\
    \bottomrule
    \end{tabular}
\end{table}
\clearpage

\subsection{Het belang van tijd in anomaliedetectie}%
\label{subsec:belang-tijd}
Tijd speelt een centrale rol in anomaliedetectie voor web- en api verkeer. Aanvallen zoals DDoS, credential stuffing en API-misbruik zullen zelden te detecteren zijn op basis van één enkele request. Deze aanvallen zullen eerder pas zichtbaar zijn voor een model wanneer ze worden geanalyseerd op een bepaalde tijdsspanne\autocite{Schummer2024}.%  Machine Learning-Based Network Anomaly Detection https://www.researchgate.net/publication/387151644_Machine_Learning-Based_Network_Anomaly_Detection_Design_Implementation_and_Evaluation
Om gedragsmatige patronen te kunnen detecteren wordt ruwe logdata samengevoegd over tijdsvensters. Dit is wanneer requests van eenzelfde client, endpoint of sessie samengebracht worden binnen een interval, dit kan bijvoorbeeld tussen de 1 en 5 minuten zijn. Per tijdsvenster wordt vervolgens de relevante features geanalysseerd\autocite{Sowmya2023}. % https://www.researchgate.net/publication/372514849_API_Traffic_Anomaly_Detection_in_Microservice_Architecture

De keuze van tijdsvenster is afhankelijk van de data waarop het ML model gebruikt zal worden, doorgaans zijn er twee mogelijkheden: tumbling windows en sliding windows. Bij tumbling windows worden er gebruik gemaakt van niet overlappende intervallen. Bij sliding windows is er sprake van een continue verschuiving van het interval om zo een gelijkmatiger beeld te geven van de gedragsverandering. Voor real-time detectie is sliding window aanbevolen doordat deze sneller reageren op plotse gedragswijzigingen\autocite{Schummer2024}. % Machine Learning-Based Network Anomaly Detection https://www.researchgate.net/publication/387151644_Machine_Learning-Based_Network_Anomaly_Detection_Design_Implementation_and_Evaluation

\subsection{Concept drift in anomaliedetectie}%
\label{subsec:concept-drift}

Machine-learningmodellen in de context van anomaliedetectie worden getrained op data die het normale verkeersgedrag van op dat moment representeert. In werkelijkheid verandert dit normaal gedrag continu: het lanceren van nieuwe features, het stijgen of dalen van verkeer van gebruikers en aanvallers die zich aanpassen. Deze verschuiving van eigenschappen en gedrag in data wordt ook wel concept drift genoemd. Wanneer concept drift optreed zonder dat het model wordt aangepast, ontstaat er een groeiende kloof tussen wat het model als normaal beschouwt en wat in de werkelijkheid op dat moment normaal is, het gevolg hiervan is een daling in detectieperformantie.   %https://dl.acm.org/doi/epdf/10.1145/3152494.3152501


\section{Feature selectie en engineering voor anomaliedetectie}
\label{sec:ML-methoden-anomalie}
\subsection{Belang van data quality}
\label{subsec:toegevoegde-waarde-ML}
% Uitleggen waarom goede data data essentieel is voor ML success 
Anomaliedetectie binnen Web- en API-beveiliging is sterk afhankelijk van de beschikbare data. Slechte data kwaliteit zoals onvolledige, inconsistente API of web metrics zorgen voor onnauwkeurige resultaten en verhogen het risico op false positives en negatives. Het verzekeren van data integriteit door correcte collectie, opslag en preprocessing is cruciaal voor effectieve anomalie detectie \Autocite{Meegle2025}.

\subsection{Van ruwe logdata naar gedragsprofielen}
\label{subsec:log-gedrags}
% Lijst de belangrijkste features voor anomaly detection in Web en API-context
Ruwe logvelden zoals IP-adres, statuscode of referrer vertellen op zichzelf weinig of een client zich al dan wel of niet normaal gedraagt. Een enkele request dat een 404 statuscode heeft is op zich niet verdacht, maar tientallen 404 requests van hetzelfde ip binnen een tijdsframe van enkele minuten kan wijzen op scanning. 
De meerwaarde voor een machine learning model ontstaat pas wanneer individuele logvelden worden samengevoegd en afgeleide kenmerken worden berekend die een gedragspatroon kunnen beschrijven over een tijdsvenster.
Dit proces wordt beschreven als feature engineering en is een belangrijke stap in anomalie detectie, zonder betekenisvolle features is het voor een goed presterend algoritme onmogelijk om een onderscheid te kunnen maken tussen normaal of afwijkend verkeer.%chandola en sowmya
Aangezien web- en API-verkeer verschillende kenmerken vertonen en andere aanvalspatronen kennen, spiltsen we binnen dit onderzoek de relevante features op in twee groepen. 
\subsubsection{API-specifieke features}
Voor het anomaliedetectie model voor API-verkeer, wordt er vertrokken van een set ruwe logvelden die rechtstreekse relevant zijn aan de API-endpoints: client IP-adres, API-endpoint, response data size, HTTP-statuscode, timestamp en request type. Deze kenmerken worden niet rechtstreeks gebruikt als input voor het model maar, dienen eerst te worden geaggregeerd over een tijdsvenster en dit expliciet per client-IP~\autocite{Sowmya2023}. 

De afgeleide features worden onderverdeeld in drie categorieën: volumetrische kenmerken, responsekenmerken en endpoint gedragskenmerken. 

\textbf{Volumetrische kenmerken} beschrijven de intensiteit waarop een client de API endpoints aanspreekt. Het totale aantal API-calls per IP binnen het tijdsvenster is de meest directe aanwijzing dat een IP mogelijks bezig is met flooding van het endpoint of geautomatiseerd misbruik \autocite{Sowmya2023}. Wanneer een IP binnen een kort tijdsvenster een afwijkend groot datavolume ontvangt, kan dit ook wijzen op het systematisch ophalen van gevoelige gegevens via de API \autocite{Sowmya2023, OWASP2023}.

\textbf{Responsekenmerken} deze kenmerken beschrijven het resultaat van uitgevoerde requests richting de API. Een verhoogd aantal 4xx-statuscodes kunnen een aanwijzing zijn van verkennend gedrag waarbij een IP systematisch endpoints bevraagt\autocite{OWASP2023}.

\textbf{Endpoint gedragskenmerken} beschrijven hoe een IP de endpoints benadert. Hierbij zijn de kernfeatures het aantal aangesproken API-endpoints en het totale aantal requests. Een afwijkend volume kan wijzen op endpoint abuse of brute force-aanvallen\autocite{Sowmya2023}. Een hoge verhouding met veel unieke endpoints in een tijdsvenster kan wijzen op het systematisch in kaart brengen van het API-oppervlakte \autocite{OWASP2023}.

\subsubsection{Web-specifieke features}
Voor het model dat anomalieën detecteert in HTTP web-verkeer, wordt vertrokken vanaf basis webserverlogvelden zoals: source IP, timestamp, method, URI, referrer. De gegevens worden gegroepeerd per source IP binnen een tijdsvenster, hierna worden voor dit specifieke ip relevante afgeleide features berekend die het gedrag van de client kunnen representeren~\autocite{Chua2024}. 

De afgeleide features worden onderverdeeld in vier categorieën: frequentiekenmerken, inhoudskenmerken, timing-kenmerken en gedragskenmerken. Deze categorisatie sluit aan bij de aanbevelingen uit de literatuur om zowel intensiteit als aard van het clientgedrag te onderzoeken \autocite{Chua2024, Schummer2024}.

\textbf{Frequentiekenmerken} beschrijven hoe intensief en op welke manier een client het systeem gebruikt binnen een bepaald tijdsvenster. Deze kenmerken worden berekend per IP-adres en geven inzicht in het algemene gebruikspatroon van een client.
De belangrijkste kenmerken zijn het totale aantal requests per IP en de verdeling van HTTP-statuscodes. Een afwijkend hoge frequentie aan requests kan wijzen op scraping of een DDoS-achtig patroon, terwijl hoge 4xx-responses kan duiden op scanning of brute force-pogingen\autocite{Chua2024, OWASP2023}.

\textbf{Inhoudskenmerken} geven een beeld over de requests zelf. Afwijkende requests bevatten vaak een afwijkend patroon zoals lange URI's die kunnen wijzen op injection-pogingen. Relevante kenmerken kunnen de gemiddelde en maximale URI-lengte zijn en de frequentie van specifieke URI-patronen binnen één IP\autocite{Chua2024}. 

\textbf{Timingkenmerken} beschrijven het tijds patroon van de requests. Normaal verkeer zal doorgaans variatie vertonen in de intervallen tussen requests, terwijl scripts en geautomatiseerde tools eerder uniforme intervallen vertonen. Daarom zijn de gemiddelde tijd en de standaardafwijking indicatieve features voor het onderscheid tussen menselijk of geautomatiseerd \autocite{Schummer2024, Benova2024}. 

\textbf{Gedragskenmerken} representeren de consistentie en diversiteit in het gedrag van de client. Voorbeelden hiervan zijn het gebruik van verschillende HTTP-methoden, verhouding tussen unieke paden, totale aantal requests en consistentie in de user-agent. Normaal verkeer zal zich doorgaans beperken tot Get en POST methoden en zal zelden in een tijdsvenster wisselen van user-agent \autocite{Benova2024, Chua2024}.  

\section{Evaluatie en performance metrics}
\label{sec:ML-methoden-anomalie}
\subsection{Evaluatie bij supervised anomaliedetectie}
\label{subsec:toegevoegde-waarde-ML}
Bij supervised anomaliedetectie beschikt de onderzoeker over ten minste één volledige gelabelde dataset waarbij elk punt expliciet geclassificeerd is als normaal of abnormaal. Dit maakt het mogelijk om de voorspellingen van het model rechtstreeks te vergelijken met de bekende grondwaarde en klassieke classificatiemetrics te berekenen \autocite{Goldstein2016}.

Precision meet het aandeel correct gedetecteerde anomalieën ten opzichte van alle als anomalie geclassificeerde observaties. Recall meet het aantal correcte anomalieën ten opzichte van het totale aantal dat zich in de dataset bevindt. Beide metrics zijn belangrijk in de evaluatie van anomaliedetectie aangezien beide zich focussen op de minderheidsklasse binnen het model (de anomalieën) en zo een nauwkeuriger beeld bieden over de detectiekwaliteit. De F1-score combineert deze twee metrics als gemiddelde, waardoor de F1-score enkel kan stijgen indien beide waarden stijgen\autocite{Chadha2020}.  

\subsection{Evaluatie bij unsupervised anomaliedetectie}
\label{subsec:toegevoegde-waarde-ML}
% Definiëren van de metrics die je gaat gebruiken voor de evaluatie van het model. 
Bij unsuperised anomaliedetectie is een betrouwbare gelabelde dataset of grondwaarde doorgaans niet beschikbaar, dit is vaak ook de reden waarom in eerste instantie gekozen is voor de implementatie van een unsupervised model. Dit heeft echter wel rechtstreeks gevolgen op de evaluatie van het model, metrics zoals recall of F1-score vereisen kennis van het totale aantal werkelijke anomalieën die zich in de dataset bevinden \autocite{Goldstein2016}.

Om toch de detectiekwaliteit van deze modellen te kunnen beoordelen worden voor deze modellen vaak alternatieve strategieën gebruikt. Een veelgebruikte strategie is handmatige validatie, waarbij de detecties gevonden door een model worden beoordeeld door een domeinexpert of door de onderzoeker zelf \textcite{Chua2024} pasten deze aanpak toe door de testdata van hun Isolation Forest-model handmatig te labelen, waarna precision, recall en F1-score berekend konden worden.

\textcite{Rochner2023} gebruikten een vergelijkbare strategie bij unsupervised anomaliedetectie in medische dossiers, waarbij domeinexperts een steekproef van de detecties beoordeelden om de precision van het model te kunnen bepalen. 

In de praktijk worden meerdere technieken vaak gecombineerd om vanuit meerdere invalshoeken een beeld te creeëren over hoe het model detecteert, de keuze hiervan is vaak afhankelijk van de beschikbare methoden. 

\section{Conclusie literatuurstudie}
\label{sec:ML-methoden-anomalie}

De literatuurstudie toont aan dat traditionele detectiesystemen zoals WAF's en signature based IDS-systemen effectief zijn in het herkennen van eerder geziene en gedocumenteerde dreigingen, maar beperkingen kunnen vertonen bij het detecteren van ongeziene of subtiele aanvalsvormen. Toenemende automatisering en complexiteit van aanvallen versterkt deze problematiek. 

Uit de vergelijking van machine learning-benaderingen blijkt dat unsupervised methoden het meest geschikt zijn wanneer gelabelde trainingsdata ontbreekt, wat in de meeste operationele security-omgevingen het geval is. 

Isolation Forest onderscheidt zich binnen deze categorie door zijn lineaire tijdcomplexiteit, lage geheugenverbruik en de eigenschap dat het anomalieën rechtstreeks isoleert in plaats van eerst normaal gedrag te modelleren. Dit maakt het bijzonder geschikt voor toepassing binnen een real-time detectiesysteem dat grote hoeveelheden logdata moet verwerken.

Daarnaast blijkt feature engineering een bepalende invloed te hebben op de detectiekwaliteit van het model. Ruwe logvelden zeggen weinig op zich, de meerwaarde ontstaat door deze velden te aggregeren over een tijdsvenster en per client-IP, hierdoor worden gedragspatronen van gebruikers zichtbaar en kunnen afwijkingen worden gedetecteerd. 

Deze inzichten vormen de basis voor de modelkeuze, feature engineering en architectuur van het proof-of-concept dat in de volgende hoofdstukken wordt uitgewerkt.
% Voorbeeldtekst van het sjabloon  




%%=============================================================================
%% Methodologie
%%=============================================================================

\chapter{\IfLanguageName{dutch}{Methodologie}{Methodology}}%
\label{ch:methodologie}

%% TODO: In dit hoofstuk geef je een korte toelichting over hoe je te werk bent
%% gegaan. Verdeel je onderzoek in grote fasen, en licht in elke fase toe wat
%% de doelstelling was, welke deliverables daar uit gekomen zijn, en welke
%% onderzoeksmethoden je daarbij toegepast hebt. Verantwoord waarom je
%% op deze manier te werk gegaan bent.
%% 
%% Voorbeelden van zulke fasen zijn: literatuurstudie, opstellen van een
%% requirements-analyse, opstellen long-list (bij vergelijkende studie),
%% selectie van geschikte tools (bij vergelijkende studie, "short-list"),
%% opzetten testopstelling/PoC, uitvoeren testen en verzamelen
%% van resultaten, analyse van resultaten, ...
%%
%% !!!!! LET OP !!!!!
%%
%% Het is uitdrukkelijk NIET de bedoeling dat je het grootste deel van de corpus
%% van je bachelorproef in dit hoofstuk verwerkt! Dit hoofdstuk is eerder een
%% kort overzicht van je plan van aanpak.
%%
%% Maak voor elke fase (behalve het literatuuronderzoek) een NIEUW HOOFDSTUK aan
%% en geef het een gepaste titel.



\subsection{Requirementsanalyse}
\label{subsec:Requirementsanalyse}
%Het doel van dit onderdeel is om in een interview met het evolane team de functionele en niet-functionele requirements zijn voor het anomaliedetectiesysteem. 

De doelstelling van dit proof of concept is het ontwikkelen van een prototype voor anomaliedetectie dat automatisch afwijkend gedrag in webverkeer identificeert op basis van Akamai access logs. Het systeem dient als aanvulling op de bestaande regelgebaseerde detectie van Akamai App \& API Protector en moet integreerbaar zijn met de bestaande monitoring-omgeving van Evolane.

\subsubsection{Functionele vereisten}
 
\begin{itemize}
    \item \textbf{Detectie van afwijkend verkeer:} het systeem moet in staat zijn om anomalieën in webverkeer te identificeren op basis van gedragspatronen, zonder afhankelijk te zijn van vooraf gedefinieerde aanvalssignaturen. Dit omvat onder andere DDoS- en DoS-patronen, botverkeer, vulnerability scanning en brute force pogingen.
    \item \textbf{Near-real-time verwerking:} het systeem moet nieuwe logs binnen een aanvaardbare vertraging kunnen verwerken en beoordelen, zodat het SOC-team tijdig kan reageren op gedetecteerde anomalieën.
    \item \textbf{Bruikbare alerts:} elke gedetecteerde anomalie moet voldoende context bevatten zodat een security-analist in staat is om het gedrag te beoordelen. 
    \item \textbf{Beperking van herhaalde alerts:} het systeem moet voorkomen dat hetzelfde IP-adres herhaaldelijk dezelfde anomalie genereert, zodat het overzicht voor het SOC-team bewaard blijft.
    \item \textbf{Vergelijking met Akamai:} het systeem moet inzicht bieden in de overlap en complementariteit met de bestaande Akamai-detectie, zodat duidelijk is welke anomalieën een meerwaarde vormen ten opzichte van de regelgebaseerde aanpak.
\end{itemize}
 
\subsubsection{Niet-functionele vereisten}
 
\begin{itemize}
    \item \textbf{Integratie met bestaande tooling:} het systeem moet integreerbaar zijn met de bestaande monitoring-omgeving van Evolane, specifiek met Grafana voor visualisatie en Elasticsearch als databron.
    \item \textbf{Beperking van false positives:} het systeem moet een aanvaardbaar lage false positive rate bereiken om alert fatigue bij het SOC-team te voorkomen. Een false positive is hierbij gedefinieerd als een normale menselijke gebruiker die onterecht als anomalie wordt geclassificeerd.
    \item \textbf{Configureerbaarheid:} de parameters van het systeem, zoals het polling-interval, de venstergrootte, de filterregels en de verbindingsgegevens, moeten centraal configureerbaar zijn zodat het systeem inzetbaar is voor meerdere klantomgevingen zonder aanpassing van de broncode.
\end{itemize}

De eigen website van Evolane werd aangeboden als testomgeving voor het proof of concept, aangezien deze regelmatig het doelwit is van aanvallen en daarmee representatief verkeer bevat voor de evaluatie van het anomaliedetectiesysteem.
%Focus op real-time detectie
%Duidelijkheid bieden aan een soc analist waarom iets gevlagd is als een anomalie
%Overzichtelijk zijn wat er precies gaande is. 

\subsection{Dataverzameling en normalisatie}
\label{subsec:dataverzameling en normalisatie}
%Het doel van dit onderdeel is om een datastroom op te zetten van de Akamai access logs zodat we python kunnen gebruiken voor deze op te schonen en het doel is om een gestructureerde set op te bouwen in CSV formaat. 

Om een reproduceerbare en herbruikbare dataset op te bouwen, wordt een datastream opgezet op basis van Akamai access logs. Deze logs bevatten onder meer informatie over client-IP-adressen, HTTP-methodes, statuscodes, responsegroottes en tijdstempels.

De logdata wordt via een geconfigureerde datastream doorgestuurd naar een Elasticsearch-cluster dat door Evolane wordt beheerd.

Dit omvat onder andere het verwijderen van onvolledige of corrupte records, het uniformeren van tijdsformaten, het omzetten van categorische variabelen, het filteren van monitoring- en health-check-verkeer en het structureren van de data in een consistent schema.

Na normalisatie worden de gegevens opgeslagen in een gestructureerd CSV-formaat. Deze gestructureerde dataset vormt de basis voor verdere analyse, feature engineering en modeltraining binnen het proof-of-concept.

\subsection{Feature engineering en gedragsmodellering}
\label{subsec:Feature engineering}
% Het doel is om op basis van de literatuurstudie een selectie maken van de kenmerken waarop het model zal draaien. Uit deze stap moet een feature set komen die geschikt is voor toepassing. 

In deze fase worden betekenisvolle kenmerken (features) geëxtraheerd uit de genormaliseerde dataset. De selectie van deze features is gebaseerd op inzichten uit de literatuurstudie rond anomaliedetectie in webverkeer.

Voorbeelden van potentiële features zijn onder meer:
- Aantal requests per client binnen een bepaald tijdsvenster  
- Verdeling van HTTP-statuscodes  
- Frequentie van specifieke API-endpoints  
- Gemiddelde en afwijkende responsegroottes  
- Tijd tussen opeenvolgende requests  

De data wordt geaggregeerd over overlappende tijdsvensters van vijf minuten met een verschuivingsstap van één minuut, gegroepeerd per client-IP.Per venster worden frequentiekenmerken berekend zoals het aantal requests en de verdeling van statuscodes, timing-kenmerken zoals de gemiddelde tijd tussen requests, inhoudskenmerken zoals URI-lengte en responsgrootte, en gedragskenmerken zoals user-agent-consistentie en de verhouding tussen unieke paden en het totale aantal requests. Op die manier worden gedragsprofielen opgebouwd die normaal gebruik representeren. 

Het resultaat van deze fase is een feature-set die geschikt is voor toepassing binnen machine-learningmodellen voor anomaliedetectie.


\subsection{Selectie machine-learning model}
\label{subsec:traditioneel-beperkingen}
% Binnen de literatuurstudie onderzoeken we welke machine learning benaderingen er effectief zijn, en hier gaan we dan een selectie maken welke geschikt zijn voor het proof of concept op basis van de beschikbare data

Op basis van de literatuurstudie en de kenmerken van de beschikbare data wordt een geschikte machine-learningbenadering geselecteerd. Aangezien de dataset voornamelijk niet-gelabelde data bevat, ligt de focus op unsupervised anomaliedetectietechnieken.

Mogelijke modellen die onderzocht worden zijn onder andere Isolation Forest, Local Outlier Factor (LOF), One-Class SVM en autoencoders. De selectie gebeurt op basis van volgende criteria:
- Geschiktheid voor niet-gelabelde data  
- Vermogen om gedragsmatige anomalieën te detecteren  
- Interpretatie van resultaten  
- Computationele haalbaarheid binnen de scope van een proof-of-concept  

Na een vergelijkende analyse wordt één primair model geselecteerd voor implementatie binnen het proof-of-concept. Indien haalbaar kan een tweede model worden gebruikt ter vergelijking van detectieperformantie.


\subsection{proof-of-concept}
\label{subsec:Poc}
% Implementeren van de gekozen ML modellen uit de vorige stap en het evalueren van de performantie en het aantal false positives. En onderzoeken hoe het model optimaal geimplementeerd kan worden met het security platform van Evolane dus doorsturen van alerts naar slack en het ontwikkelen van een eenvoudig dashboard voor visualisatie. 
Om de theoretische inzichten te toetsen wordt er een Proof of Concept uitgewerkt waarin een anomaliedetectiemethode wordt geïmplementeerd. De detectieperformantie wordt geëvalueerd aan de hand van precision en false positive ratio. Precision geeft aan welk aandeel van de als anomalie geclassificeerde datapunten daadwerkelijk afwijkend gedrag vertoont. De false positive ratio is essentieel voor het bepalen in welke mate het prototype het aantal foutieve detecties kan beperken. Daarnaast wordt een handmatige validatie uitgevoerd door een security-analist om de kwaliteit van de detecties inhoudelijk te beoordelen.
Naast de technische evaluatie wordt ook de operationele integratie onderzocht. Hierbij wordt nagegaan hoe het prototype gekoppeld kan worden aan de bestaande monitoring van Evolane. De detectieresultaten worden ontsloten via een API en gevisualiseerd in een Grafana-dashboard, zodat het SOC-team de gedetecteerde anomalieën kan inspecteren.

\subsection{Conclusie en advies}
\label{subsec:conclusie en advies}
% Analyse opbouwen in welke mate het machine learning gebasseerd anomaliedetectiesysteem een meerwaarde is binnen de security omgeving van Evolane en het formuleren van advies. 
Tot slot wordt er een conclusie geformuleerd op basis van de bekomen resultaten. Hierbij wordt nagegaan in welke mate een machine-learning gedreven aanpak een meerwaarde biedt ten opzicht van klassieke detectie mechanismen. Op basis van de analyse wordt een onderbouwd advies geformuleerd met aandacht voor toepasbaarheid, voordelen, beperkingen en mogelijke vervolgstappen.





\chapter{\IfLanguageName{dutch}{Proof of Concept}{Proof of Concept}} 
\label{ch:proof-of-concept}

Dit hoofdstuk beschrijft de volledige opzet en implementatie van het proof-of-concept voor het machine learning-gedreven anomaliedetectiesysteem binnen de omgeving van evolane specifiek voor web verkeer. 
Hier wordt opeenvolgend de context en dataomgeving van Evolane toegelicht, dataverzameling en normalisatie besproken, feature engineering uitgewerkt, de modelkeuze onderbouwd, het trainingsproces beschreven en tot slot de operationele implementatie binnen de bestaande omgeving van Evolane geschetst.

\section{Context, omgeving en data}
\label{sec:context-omgeving-data}


\subsection{De Evolane-omgeving}
\label{subsec:evolane-omgeving}

Evolane beveiligt het web- verkeer van haar klanten via het Akamai App & API protector platform. Inkomende HTTP-requesten komen binnen op de Akamai edge servers waarna ze worden geëvalueerd op basis van regelgebasseerde beveiligingscontroles. Elke request genereert een gedetailleerd access log dat informatie bevat over het request. Deze access logs bevatten onder andere het client-IP-adres, opgevraagde pad, de HTTP-methode, dee statuscode, de responsegrootte, de user-agent, geolocatie en eventueel een Akamai beveiligingsregel indien deze geraakt zou zijn. 

De testomgeving die door Evolane ter beschikking werd gesteld voor dit proof of concept bevat uitsluitend webverkeer. Om die reden ligt de focus binnen dit proof of concept op de detectie van afwijkend gedrag binnen webverkeer. Daarnaast is het ook de doelstelling om binnen de beschikbare tijd van deze bachelorproef een degelijk werkbaar proof of concept te realiseren.
%oplijsting maken rond de features in de logs kijken of een tabel dit niet beter illustreert
\subsection{De datastream naar Elasticsearch}
\label{subsec:datastream-elasticsearch}

De Akamai access logs worden via een geconfigureerde datastream doorgestuurd naar een Elasticsearch-cluster dat door Evolane wordt beheerd. Deze datastream schrijft de logs weg naar een index in de cluster met het patroon \textit{datastream2-*}. Een enkel HTTP-request binnen deze index bevat de volledige set aan velden die Akamai per request genereert/opslaat. De belangrijkste velden binnen dit onderzoek zijn de volgende: 

\begin{itemize}
    \item \textbf{reqTimeSec}: het tijdstip van de request als Unix-epoch timestamp in milliseconden
    \item \textbf{cliIP}: het IP-adres van de client
    \item \textbf{reqPath}: het opgevraagde URI-pad
    \item \textbf{reqMethod}: de gebruikte HTTP-methode (GET, POST, PUT, DELETE, PATCH, HEAD, OPTIONS, TRACE )
    \item \textbf{statusCode}: de HTTP-statuscode van het antwoord
    \item \textbf{bytes} en \textbf{totalBytes}: de grootte van de response in bytes
    \item \textbf{UA} en \textbf{user\_agent}: de user-agent string en een user-agent object met informatie over het besturingssysteem, browser en het apparaat
    \item \textbf{country}, \textbf{city}, \textbf{state}: geolocatiegegevens van de client
    \item \textbf{referer}: de referrer-header van de request
    \item \textbf{queryStr}: de query string parameters
    \item \textbf{reqHost} en \textbf{reqPort}: de host en poort waarop de request binnenkwam
    \item \textbf{tlsVersion}: de TLS-versie
    \item \textbf{cacheStatus}: of de response uit de cache werd geserveerd
    \item \textbf{transferTimeMSec} en \textbf{turnAroundTimeMSec}: timinggegevens van de request
    \item \textbf{rspContentType}: het content type van het antwoord
    \item \textbf{aapRule}: een gestructureerd veld dat aangeeft welke Akamai-beveiligingsregels zijn getriggerd, inclusief de regelcategorie en het bijhorende beleid
\end{itemize}

Het \textit{aapRule}-veld is relevant voor dit onderzoek. Wanneer de velden \textit{rules} en \textit{category} leeg zijn, wilt dit zeggen dat Akamai geen beveiligingsregel heeft gevonden voor dit specifieke request. Dit is het verkeer dat we specifiek willen observeren of we hier anomaliën in kunnen vinden. Requests die door Akamai APP & API protector niet als verdacht zijn aangemerkt, maar mogelijk wel afwijkend gedrag vertonen. 


%Training model uitleggen 
\subsection{Trainingsdata}
\label{subsec:trainingsdata}

Voor het trainen van het model is een dataset opgehaald uit de elasticsearch-cluster over de periode van 18 maart 2026 tot en met 2 april 2026. Dit is zo gekozen aangezien dit het hoogst mogelijke aantal access logs waren die in één keer opgehaald konden worden, door dat de elasticsearch de logs maar 14 dagen bewaard. 
Deze dataset is voldoende representatief voor het normale verkeer van de Evolane website te kunnen beschrijven aangezien het zowel weekdagen als weekenddagen omvat waar het verkeer mogelijks wat kan verschillen. De trainigsdataset bevat alle requests uit deze periode inclusief de requests waarvoor Akamai wel of geen beveiligingsregel heeft getriggerd. De Akamai regelvelden \textit{rules} en \textit{category} zullen niet gebruikt worden als effectieve input voor het model, maar dienen achteraf ter referentie om de vergelijking te kunnen maken met wat het model detecteerde en wat Akamai detecteerde.

\section{Dataverzameling en normalisatie} 
\label{sec:dataverzameling-normalisatie}

\subsection{Ophalen van de logdata} 
\label{subsec:ophalen-logdata}

Het ophalen van de logs gebeurt op twee verschillende manieren, afhankelijk van het doel. Continu voor de real-time detectiepipeline of éénmalig voor het samenstellen van de trainingsdata. 

\subsubsection{Trainingsdata} 
\label{subsubsec:ophalen-trainingsdata}

Voor de trainingsdata van het model wordt er eenmalig een dataset opgehaald uit de elasticsearch van Evolane. Deze wordt vervolgens opgeslagen in een CSV bestand zodat we hier later het model op kunnen trainen. Dit CSV-bestand bevat alle beschikbare Akamai-logvelden per request die opgeslagen zijn in de elasticsearch. 
Het trainingsscript is instaat dit CSV bestand in te lezen, het uitvoeren van preproccesing en feature engineering en tot slot het model te trainen op de data. Doordat de trainingsdata wordt opgeslagen in een CSV is het mogelijk om het model hierop opnieuw te trainen zonder een verbinding te moeten maken met elasticsearch. 

\subsubsection{Real-time pipeline} 
\label{subsubsec:ophalen-pipeline}

De operationele pipeline (\textit{anomaly\_pipeline.py}) is geconfigureerd om logdata continu rechtstreeks uit de elasticsearch op te halen. Standaard zal het elke 30 seconden (configureerbaar in \textit{config.yaml}) opnieuw een poll uitvoeren naar nieuwe logs die binnengekomen zijn in elasticsearch.
Hierbij wordt er gebruik gemaakt van de \textit{search\_after}-API van Elasticsearch voor paginatie: de pipeline onthoudt wat de laatste query was die hij op elasticsearch succesvol heeft uitgevoerd en vraagt bij de volgende poll enkel een recentere page op, dit wordt zo gedaan om te voorkomen dat herhaaldelijk dezelfde data zal worden opgehaald. 

De query voor het ophalen van logs maakt gebruik van het \textit{reqTimeSec}-veld en zal enkel requests ophalen die nieuwer zijn dan de timestamp in de buffer van de pipeline. Intern wordt het tijdsveld bijgehouden als Unix-epoch in seconden en bij de query wordt deze terug omgezet naar milliseconden het formaat dat elasticsearch verwacht. 

In de pipeline wordt er expliciet gespecifieerd welke bronvelden uit elasticsearch opgeslagen worden, dit is om het geheugengebruik van de pipeline zo goed mogelijk te beheren. 

\subsection{Preprocessing en normalisatie} 
\label{subsec:preprocessing-normalisatie}

Voor de ruwe logdata geschikt is voor feature engineering ondergaat deze een reeks preprocessing-stappen die geïmplementeerd zijn in de \textit{preprocess\_logs}-functie (voor de real-time pipeline) en de \textit{preprocess\_csv}-functie (voor het trainingsscript):

\textbf{Tijdstempelnormalisatie.} Het veld \textit{reqTimeSec} kan afhankelijk van de Elasticsearch-configuratie binnenkomen als seconden of als milliseconden. Het script detecteert automatisch welk formaat wordt gebruikt door de mediaan te controleren: indien deze groter is dan $10^{12}$ wordt aangenomen dat de waarden in milliseconden staan en worden ze gedeeld door 1000. Vervolgens wordt \textit{reqTimeSec} omgezet naar leesbare UTC-timestamp. 

{\small
\begin{lstlisting}[style=custompython,caption={Detecteren formaat van timestamp},label={lst:dlp-selenium}, captionpos=b, basicstyle=\small\ttfamily]
df["reqTimeSec"] = pd.to_numeric(df["reqTimeSec"], errors="coerce")
df.dropna(subset=["reqTimeSec"], inplace=True)

if df["reqTimeSec"].median() > 1e12:
    logging.info("reqTimeSec in milliseconden, omzetten naar seconden")
    df["reqTimeSec"] = df["reqTimeSec"] / 1000.0

df["timestamp"] = pd.to_datetime(df["reqTimeSec"], unit="s", utc=True)
\end{lstlisting}
}

\textbf{Numerieke velden.} Velden zoals \textit{statusCode}, \textit{bytes}, \textit{objSize}, \textit{transferTimeMSec}, \textit{turnAroundTimeMSec},\textit{totalBytes} en \textit{rspContentLen} worden gecast naar numerieke waarden zodat deze door het model verwerkt kunnen worden. Ontbrekende/ongeldige waarden worden eerst omgezet naar NaN en vervolgens tot nul.

\textbf{User-agent ontleed.} Het \textit{user\_agent}-veld bevat een dictionary-achtige string met informatie over het besturingssysteem, de browser en het apparaat gebruikt door de client. Deze wordt ontleed in drie afzonderlijke kolommen: \textit{os\_name}, \textit{browser\_name} en \textit{device\_name}. Bij ongeldige of ontbrekende waarden worden standaardwaarde: (\enquote{Unknown}) toegekend aan de features. Dit wordt uitgevoerd door de functie: \textit{parse_user_agent_field}.

\textbf{Akamai-regel ontleed.} Het \textit{aapRule}-veld wordt op vergelijkbare wijze ontleed als de User-agent. Hieruit worden vier afgeleide velden berekend: een lijst van getriggerde regels (\textit{rules\_list}), een boolean die aangeeft of er een regel is getriggerd (\textit{has\_rule}), een boolean die aangeeft of het om een bot-gerelateerde regel gaat (\textit{is\_bot\_rule}) en het totaal aantal getriggerde regels (\textit{rule\_count}). Deze waarden worden ter correlatie gebruikt en niet als features voor het model. 

\textbf{Filtering van monitoring- en SLA-verkeer.} Bepaalde user-agents en request-paden worden uitgefilterd om te voorkomen dat monitoring- en health-check-verkeer het model beïnvloedt. Het \textit{config.yaml} specificeert welke user-agents (zoals \textit{Akamai\_Site\_Analyzer} en \textit{updown.io}) en paden (zoals \textit{sla-test.html}, \textit{health} en \textit{healthcheck}) moeten worden uitgesloten door de pipeline.

\textbf{Afgeleide velden.} Tot slot worden bijkomende velden berekend die later dienen als basis voor de feature engineering: de lengte van het URI-pad (\textit{uri\_length}), het request een query string bevat (\textit{has\_query\_string}), of er een referrer aanwezig is (\textit{has\_referer}), de statuscodeklasse als variabele met een categorie (\textit{status\_class}: 2xx, 3xx, 4xx of 5xx) en tot slot of het request via HTTPS binnenkwam (\textit{is\_https}).



Na de preprocessing van de data wordt deze nog gesorteerd op wanneer de requests binnen kwamen. 

\section{Feature engineering} 
\label{sec:feature-engineering}

Zoals beschreven in de literatuurstudie (sectie~\ref{sec:belang-tijd} en~\ref{sec:feature-selectie}) is het belangrijk om ruwe logdata te aggregeren over tijdsvensters en per client-IP, zodat gedragspatronen op deze manier zichtbaar worden. Het proof-of-concept implementeert een sliding window-benadering voor de feature engineering.

\subsection{Sliding window-configuratie} 
\label{subsec:sliding-window-configuratie}

De feature engineering maakt gebruik van overlappende tijdsvensters met de volgende configuratie:

\begin{itemize}
    \item \textbf{\textit{Venstergrootte}}: 300 seconden (5 minuten)
    \item \textbf{\textit{Verschuivingsstap}}: 60 seconden (1 minuut)
    \item \textbf{\textit{Minimum aantal requests}}: 3 requests per venster per client-IP
\end{itemize}

De keuze voor een venstergrootte van vijf minuten is gebaseerd op de aanbevelingen uit de literatuur~\autocite{Schummer2024} en biedt een balans tussen het detecteren van kort durende aanvallen en het opvangen van voldoende informatie om het gedrag te kunnen bepalen. De verschuivingsstap van één minuut zorgt ervoor dat het systeem snel reageert op gedragswijzigingen: elke minuut zal de pipeline een nieuw vijf-minuten venster evalueren. Binnen de pipeline worden Client IPs waar minder dan 3 logs voor gevonden zijn niet verwerkt aangezien dit onvoldoende data is om features voor te berekenen. 

Het resultaat van deze stap is een DataFrame waarbij elke rij één combinatie van client-IP en tijdsvenster voorstelt, voorzien van alle berekende features.

\subsection{Frequentie features}
\label{subsec:frequentiekenmerken}

De frequentie features beschrijven hoe intensief een client het systeem gebruikt binnen een tijdsvenster:

\begin{itemize}
    \item \textbf{\textit{request\_count}}: het totale aantal requests van het client IP-adres binnen het venster. Een abnormaal hoog aantal kan wijzen op scraping, DDoS-patronen of geautomatiseerde aanvallen.
    \item \textbf{\textit{requests\_per\_second}}: het aantal requests gedeeld door de duur van het venster (berekend als het verschil tussen de laatste en het eerste timestamp binnen het venster). Deze feature schaalt het volume ten opzichte van de effectieve duur van het venster.
    \item \textbf{\textit{unique\_paths}}: het aantal unieke URI-paden dat werd aangevraagd door het client ip. Een zeer laag aantal in combinatie met een hoog requestvolume kan wijzen op endpoint hammering.
    \item \textbf{\textit{path\_diversity\_ratio}}: de verhouding tussen het aantal unieke paden en het totale aantal requests. Een lage ratio wijst op repetitief gedrag gericht op een beperkt aantal endpoints, wat indicatief kan zijn voor brute force- of URI-misbruik.
\end{itemize}

{\small
\begin{lstlisting}[style=custompython,caption={Berekenen van frequentie features},label={lst:dlp-selenium}, captionpos=b, basicstyle=\small\ttfamily]
def compute_window_features(group: pd.DataFrame) -> dict:
    request_count = len(group)

    window_duration = (group["reqTimeSec"].max() - group["reqTimeSec"].min())
    rps = request_count / max(window_duration, 1.0)

    unique_paths = group["reqPath"].nunique()
    path_diversity = unique_paths / request_count if request_count > 0 else 0
\end{lstlisting}
}

\subsection{Statuscodedistributie} 
\label{subsec:statuscodedistributie}

De verdeling van HTTP-statuscodes biedt inzicht in het resultaat van de uitgevoerde requests:

\begin{itemize}
    \item \textbf{\textit{rate\_2xx}}: het aantal succesvolle requests. Normaal verkeer vertoont doorgaans een 2xx statuscodes.
    \item \textbf{\textit{rate\_3xx}}: het aantal redirects.
    \item \textbf{\textit{rate\_4xx}}: het aantal client-fouten (400, 403, 404, enz.). Een verhoogd aantal 4xx-responses kan wijzen op scanning, pad-enumeratie of ongeautoriseerde toegangspogingen~\autocite{OWASP2023}.
    \item \textbf{\textit{rate\_5xx}}: het aantal server-fouten. Een hoge 5xx-ratio kan duiden op pogingen tot exploitatie die serverfouten veroorzaken.
\end{itemize}

\subsection{Timing-kenmerken} 
\label{subsec:timing-kenmerken}

De timing-kenmerken beschrijven het gedrag gehanteerd door het client IP en kunnen het tijdsgebonden patroon tussen requests beschrijven:

\begin{itemize}
    \item \textbf{\textit{avg\_inter\_request\_time}}: de gemiddelde tijd tussen opeenvolgende requests. Menselijke gebruikers vertonen doorgaans variatie in intervallen, terwijl geautomatiseerde tools vaak meer uniforme intervallen vertonen.
    \item \textbf{\textit{std\_inter\_request\_time}}: de standaardafwijking van de tussentijden. Een zeer lage standaardafwijking wijst op mechanisch en dus potentieel geautomatiseerd gedrag.
    \item \textbf{\textit{min\_inter\_request\_time}}: de kortste tijd tussen twee opeenvolgende requests. Extreem korte minimale intervallen kunnen duiden op scripted requests.
\end{itemize}

{\small
\begin{lstlisting}[style=custompython,caption={Berekenen van timing features},label={lst:dlp-selenium}, captionpos=b, basicstyle=\small\ttfamily]
    if n > 1:
        inter_times = group["reqTimeSec"].diff().dropna() 
        avg_inter_time = inter_times.mean()
        std_inter_time = inter_times.std()
        min_inter_time = inter_times.min()
    else:
        avg_inter_time = 0.0
        std_inter_time = 0.0
        min_inter_time = 0.0
\end{lstlisting}
}

\subsection{Inhoudskenmerken} 
\label{subsec:inhoudskenmerken}

De inhoudskenmerken beschrijven de aard en omvang van de requests en responses:

\begin{itemize}
    \item \textbf{\textit{avg\_uri\_length}}: de gemiddelde lengte van de opgevraagde URI's. Aanvallen zoals injection-pogingen genereren vaak abnormaal lange URI's~\autocite{Chua2024}.
    \item \textbf{\textit{max\_uri\_length}}: de maximale URI-lengte binnen het venster. Een enkele extreem lange URI kan indicatief zijn voor een injection-aanval.
    \item \textbf{\textit{avg\_response\_bytes}}: de gemiddelde responsgrootte. Afwijkingen in responsgrootte kunnen wijzen op succesvolle data-exfiltratie.
    \item \textbf{\textit{total\_bytes\_out}}: het totale aantal uitgaande bytes. Een hoog totaalvolume kan wijzen op het ophalen van een abnormale hoeveelheid data.
    \item \textbf{\textit{avg\_transfer\_time\_ms}}: de gemiddelde transfertijd in milliseconden.
\end{itemize}

\subsection{Methode- en gedragskenmerken} 
\label{subsec:methode-gedragskenmerken}

\begin{itemize}
    \item \textbf{\textit{unique\_methods}}: het aantal verschillende HTTP-methoden dat werd gebruikt. Normaal webverkeer beperkt zich doorgaans tot GET en POST, terwijl aanvallers PUT, DELETE of andere methoden inzetten.
    \item \textbf{\textit{non\_get\_ratio}}: het aantal niet-GET-requests. Een hoog aantal niet-GET-requests in combinatie met andere afwijkingen kan wijzen op misbruik van de logica.
    \item \textbf{\textit{unique\_countries}}: het aantal verschillende landen van waaruit requests binnen één venster voor hetzelfde IP werden waargenomen. Hoewel dit veld doorgaans gelijk is aan 1, kan een afwijking wijzen op een anomalie.
    \item \textbf{\textit{has\_referer\_ratio}}: het aantal requests met een geldige referrer-header. Geautomatiseerde tools sturen vaak geen referrer mee, waardoor een ratio van nul verdacht kan zijn bij een hoog aantal requests.
    \item \textbf{\textit{has\_query\_ratio}}: het aantal requests met query string parameters.
    \item \textbf{\textit{https\_ratio}}: het aantal requests dat via HTTPS binnenkwam.
    \item \textbf{\textit{unique\_content\_types}}: het aantal verschillende response-content-types. Een hoge diversiteit kan wijzen op het systematisch verkennen van de applicatie.
    \item \textbf{\textit{cache\_hit\_ratio}}: het aantal requests geserveerd uit de cache.
\end{itemize}

\subsection{User-agent-consistentie}  
\label{subsec:user-agent-consistentie}

\begin{itemize}
    \item \textbf{\textit{unique\_browsers}}: het aantal verschillende browsers uit het user agent veld dat werd waargenomen voor hetzelfde IP-adres binnen het venster. Normaal gebruikersverkeer komt doorgaans van één browser; meerdere browsers binnen korte tijd kan wijzen op bot-activiteit of het gebruik van wisselende user-agents om detectie te vermijden.
    \item \textbf{\textit{unique\_os}}: het aantal verschillende besturingssystemen. Vergelijkbaar met het voorgaande: meerdere besturingssystemen voor één IP-adres is abnormaal.
\end{itemize}

\subsection{Akamai-correlatie-velden} 
\label{subsec:akamai-correlatie-velden}

Naast de features die voor het model gebruikt worden, zijn er ook drie correlatie-velden berekend die niet als input voor het model dienen, maar die bij de evaluatie worden gebruikt om de voorspellingen van het model te vergelijken met de bestaande Akamai-detectie:

\begin{itemize}
    \item \textbf{\textit{\_akamai\_has\_rule\_ratio}}: het aantal requests waarvoor Akamai een beveiligingsregel heeft gevonden.
    \item \textbf{\textit{\_akamai\_bot\_rule\_ratio}}: het aantal requests waarvoor specifiek een bot-gerelateerde regel is getriggerd.
    \item \textbf{\textit{\_akamai\_avg\_rule\_count}}: het gemiddelde aantal getriggerde regels per request.
\end{itemize}

Deze velden worden bewust uitgesloten van de modelinput om te voorkomen dat het model de Akamai-detectie repliceert in plaats van aanvullende anomalieën te ontdekken.

\section{Modelkeuze} 
\label{sec:modelkeuze}

\subsection{Ongeschiktheid van supervised benaderingen voor de beschikbare data}
\label{subsec:ongeschiktheid-supervised}

Zoals beschreven in de literatuurstudie vereisen supervised learning-algoritmen een dataset waarbij elk datapunt specifiek gelabeld is als normaal of anomaal gedrag. Binnen de akamai access logs is het \textit{aapRule}-veld het enige veld dat een indicatie geeft over de aard van het requets. Hoewel het lijkt alsof dit veld als label kan dienen is het ongeschikt als grondwaarde voor een supervised learning model. 

Ten eerste is het veld slechts een indicatie van onder andere regelgebaseerde detectie, niet een label dat werkelijk beschrijft of het gaat om een normaal of anomaal requets. Het betekent uitsluitend dat er voor dit specifieke request geen Akamai beveiligingsregel is getriggerd. 

Ten tweede dekt het \textit{aapRule}-veld niet alles van afwijkend gedrag. De Akamai-regels zijn slechts gebaseerd op bestaande aanvalssignaturen, botsignaturen en gekende patronen. Nieuwe aanvalsvormen, subtiele gedragsafwijkingen of voorheen ongekende bots worden mogelijks niet door deze regels gevlagd. Een supervised model getrained op basis van deze regels zou dezelfde blinde vlekken leren en geen meerwaarde bieden ten opzichte van het huidige systeem. 

\subsection{Onderbouwing voor Isolation Forest} 
\label{subsec:onderbouwing-isolation-forest}

Op basis van de vergelijkende analyse uit de literatuurstudie (tabel~\ref{tab:vergelijking-ml-modellen}) en de kenmerken van de beschikbare data is gekozen voor Isolation Forest als primair anomaliedetectiemodel. 

\textbf{Geschiktheid voor niet-gelabelde data.} De beschikbare trainingsdata is niet gelabeld: er zijn geen expliciete annotaties beschikbaar die aangeven welke requests of sessies daadwerkelijk kwaadaardig zijn. Isolation Forest is een unsupervised algoritme dat geen gelabelde data vereist en in plaats daarvan anomalieën detecteert door het isoleren van datapunten die afwijken van de meerderheid~\autocite{Liu2008}.

\textbf{Efficiëntie en schaalbaarheid.} Isolation Forest beschikt over een lineaire tijdcomplexiteit en een laag geheugenverbruik~\autocite{Lesouple2021}, wat het geschikt maakt voor de verwerking van grote hoeveelheden logdata. Dit is een bepalend voordeel ten opzichte van methoden zoals Local Outlier Factor, die bij grote datasets meer rekentijd vereisen vanwege de berekening van de k-nearest neighbors.

\textbf{Directe isolatie van afwijkingen.} In tegenstelling tot andere methoden die eerst een profiel opbouwen van normaal gedrag en vervolgens afwijkingen daarvan evalueren, focust Isolation Forest zich rechtstreeks op het isoleren van anomalieën. Het algoritme vertrekt vanuit de observatie dat anomalieën zeldzaam zijn en afwijkende kenmerken bezitten, waardoor ze met minder splitsingen in de decision trees geïsoleerd kunnen worden~\autocite{Liu2008}. Deze aanpak maakt het model geschikt voor het detecteren van verschillende soorten anomalieën zonder dat het model vooraf opgedragen moet worden over welke patronen afwijkend zijn.

\textbf{Interpreteerbaarheid via anomaliescores.} Isolation Forest genereert een continue anomaliescore per datapunt, wat de mogelijkheid biedt om anomalieën te rangschikken op basis van hoe ernstig het model deze scoort. Dit is waardevol voor SOC-analisten die dienen prioriteiten op te stellen rond welke aomalieën eerst verder onderzocht worden. 

\textbf{Haalbaarheid binnen de scope van een proof-of-concept.} Alternatieven zoals autoencoders vereisen een aanzienlijke mate van architectuuroptimalisatie en hyperparameter-tuning, wat minder haalbaar is binnen het tijdskader van dit onderzoek. Isolation Forest biedt daarentegen een robuuste baseline met beperkte configuratievereisten.

\subsection{Afweging ten opzichte van alternatieven} 
\label{subsec:afweging-alternatieven}

Local Outlier Factor (LOF) werd overwogen vanwege de sterke capaciteit voor het detecteren van lokale anomalieën, maar de hogere rekencomplexiteit en de gevoeligheid voor de parameterkeuze van $k$ maken het minder geschikt voor een real-time detectiesysteem die een grote hoeveelheid data dient te verwerken.

One-Class SVM werd overwogen als semi-supervised alternatief. Dit model vereist echter zorgvuldige afstemming van de kernelfunctie en de bijbehorende hyperparameters, en de trainingstijd schaalt niet-lineair met de datasetgrootte, wat de haalbaarheid binnen dit proof-of-concept beperkt.

\section{Training van het model} 
\label{sec:training-model}

\subsection{Trainingsproces} 
\label{subsec:trainingsproces}

De training van het Isolation Forest-model wordt uitgevoerd door het script \textit{train\_model.py}. Het trainingsproces verloopt als volgt:

\begin{enumerate}
    \item \textbf{Inladen van de CSV-data}: de eerder opgehaalde en opgeslagen ruwe logdata wordt ingeladen door het python script.
    \item \textbf{Preprocessing}: dezelfde preprocessing-pipeline als beschreven in sectie~\ref{subsec:preprocessing-normalisatie} wordt toegepast op de trainingsdata.
    \item \textbf{Feature engineering}: de sliding window-benadering genereert feature-vectoren over de volledige trainingsperiode.
    \item \textbf{Standaardisatie}: alle feature-waarden worden geschaald met behulp van een \textit{StandardScaler} uit scikit-learn. Van elke waarde wordt het gemiddelde van die specifieke feature afgetrokken en vervolgens gedeeld door de standaardafwijking van deze feature, dit zorgt ervoor dat elke feature tot een vergelijkbare schaal wordt gebracht. Dit is belangrijk in Isolation Forest aangezien dit algoritme gevoelig is aan aan verschillen tussen features. Wanneer één feature een veel grotere numerieke waarde heeft zoals total_bytes_out, zal deze vaker gekozen worden voor splitsingen. Dit kan ertoe leiden dat features met een kleinere range minder invloed zouden hebben op het model zoals bijvoorbeeld path_diversity_ratio. De standard scaler zorgt er dus voor dat elke feature evenveel doorweegt. 
    \item \textbf{Modeltraining}: het Isolation Forest-model wordt getraind op de geschaalde trainingsdata.
\end{enumerate}

In het trainingsscript worden waarden die oneindig zijn (\textit{inf}) of simpelweg ontbreken (\textit{NaN}) voor het stadaardisatie proces vervangen door nul om fouten tijdens de training te voorkomen.

{\small
\begin{lstlisting}[style=custompython,caption={Gebruik van standardscaler},label={lst:dlp-selenium}, captionpos=b, basicstyle=\small\ttfamily]
    def score_features(features_df: pd.DataFrame, model, scaler) -> pd.DataFrame:

    X = features_df[MODEL_FEATURES].copy()

    X.replace([np.inf, -np.inf], np.nan, inplace=True)
    X.fillna(0, inplace=True)

    X_scaled = scaler.transform(X)

    features_df = features_df.copy()
    features_df["anomaly_score"] = model.decision_function(X_scaled)
    features_df["is_anomaly"] = model.predict(X_scaled) == -1

    return features_df
\end{lstlisting}
}

\subsection{Modelparameters} 
\label{subsec:modelparameters}

Het Isolation Forest-model wordt geconfigureerd met de volgende parameters:

\begin{itemize}
    \item \textbf{\textit{n\_estimators = 200}}: het aantal decision trees in het ensemble. Een hoger aantal bomen verhoogt de stabiliteit van de anomaliescores berekend door het model doordat de gemiddelde padlengte over meer bomen wordt berekend. De waarde van 200 biedt een goede balans tussen stabiliteit en rekentijd.
    \item \textbf{\textit{contamination = 0.10}}: de verwachte proportie anomalieën in de trainingsdata. Deze parameter bepaalt de drempelwaarde waarboven een datapunt als anomalie wordt geclassificeerd. Een waarde van 10\%  is gekozen als schatting, wat betekent dat het model de 10\%  meest afwijkende datapunten als anomalie bestempelt.
    \item \textbf{\textit{random\_state = 42}}: een vaste seed voor de random number generator, wat de reproduceerbaarheid van de trainingsresultaten garandeert.
    \item \textbf{\textit{n\_jobs = -1}}: alle beschikbare CPU-cores worden ingezet voor parallelle training, wat de trainingstijd verkort.
\end{itemize}

\subsection{Model-artefacten} 
\label{subsec:model-artefacten}

Na de training worden drie componenten opgeslagen in de \textit{model/}-directory:

\begin{itemize}
    \item \textbf{\textit{isolation\_forest\_model.joblib}}: het getrainde Isolation Forest-model, geserialiseerd met joblib.
    \item \textbf{\textit{scaler.joblib}}: de gefitte StandardScaler, die noodzakelijk is om nieuwe data op dezelfde schaal te transformeren als de trainingsdata vooraleer deze door het getrainde model wordt beoordeeld.
    \item \textbf{\textit{model\_metadata.yaml}}: een metadata-bestand met alle informatie over de training van het model, waaronder het aantal trainingssamples, het aantal features, de gebruikte contamination-waarde, de score-distributie (gemiddelde, standaardafwijking, minimum en maximum) en de trainingstijd.
\end{itemize}

\subsection{Hertraining en modelonderhoud}
\label{subsec:hertraining-modelonderhoud}

Verkeerspatronen op een web omgeving zijn niet statisch: toevoeging van nieuwe applicaties, variaties die seizoensgebonden zijn, nieuwe klanten of diensten die aangeleverd kunnen worden. Dit voorval werd eerder beschreven in de literatuurstudie als (sectie~\ref{sec:concept-drift}). 

Om dit op te vangen is het proof of concept zo opgezet zodat het model eenvoudig opnieuw getrained kan worden op een nieuwe dataset. De procedure vereist enkel het ophalen van een nieuwe dataset uit de elasticsearch van evolane die minimaal twee weken aan Akamai access-logs data bevat en het opnieuw uitvoeren van het trainingsscript. De bestaande modelartefacten worden hierbij overschreven. 

% bepalen of dit er in blijft of niet aangezien de werking hiervan niet 100 procent is door de windows
\subsection{Validatie tegen Akamai} 
\label{subsec:validatie-akamai}

Als bijkomende validatiestap vergelijkt het trainingsscript de voorspellingen van het model met de Akamai-detectie via het \textit{\_akamai\_has\_rule\_ratio}-veld. Deze vergelijking levert vier categorieën op:

\begin{itemize}
    \item \textbf{Beide vlaggen}: het model en Akamai beschouwen het verkeer als verdacht.
    \item \textbf{Alleen model}: het model detecteert een anomalie die Akamai niet heeft opgepikt. Dit zijn de meest interessante gevallen, aangezien ze de potentiële meerwaarde van het ML-model aantonen.
    \item \textbf{Alleen Akamai}: Akamai heeft een regel getriggerd die het model niet als anomalie beschouwt. Dit kan voorkomen wanneer de regelgebaseerde detectie reageert op specifieke signaturen die op gedragsniveau niet afwijkend zijn.
    \item \textbf{Geen van beide}: zowel het model als Akamai beschouwen het verkeer als normaal.
\end{itemize}

\section{Implementatie in de Evolane-omgeving} 
\label{sec:implementatie-evolane}

\subsection{Architectuuroverzicht} 
\label{subsec:architectuuroverzicht}

De operationele implementatie bestaat uit een Python-applicatie (\textit{anomaly\_pipeline.py}) die als continu draaiend proces op een virtuele machine binnen de Evolane-infrastructuur wordt uitgevoerd. In elke pollcyclus haalt de pipeline nieuwe logs op uit Elasticsearch, voert het preprocessing en feature engineering uit, scoort de resulterende feature-vectoren met het getrainde model en slaat deze resultaten op in het \textit{anomalies\_detected.csv} bestand. 

\subsection{Real-time polling en rolling buffer} 
\label{subsec:polling-rolling-buffer}

De pipeline pollt elke 30 seconden (configureerbaar via \textit{config.yaml}) naar nieuwe logentries in Elasticsearch. Hierbij wordt de \textit{search\_after}-API gebruikt voor het efficiënt ophalen en dubbele logs te voorkomen, in plaats van de oudere scroll-API. Elk opgehaald log wordt aan een rolling buffer toegevoegd die de afgelopen 10 minuten (tweemaal de venstergrootte) aan data bewaart. Logs die ouder zijn dan deze buffer worden automatisch verwijderd om het geheugengebruik stabiel te houden.

\subsection{Deduplicatie en cooldown} 
\label{subsec:deduplicatie-cooldown}

Om te voorkomen dat een enkel verdacht IP-adres herhaaldelijk dezelfde anomalie genereert, implementeert de pipeline een deduplicatiemechanisme. Per pollcyclus wordt voor elk IP-adres enkel het venster met de laagste (meest anomale) score behouden. Daarnaast wordt een cooldown-mechanisme toegepast: nadat een IP-adres als anomalie is gemeld wordt het voor de duur van de venstergrootte (5 minuten) onderdrukt. Dit voorkomt dat een ip adres telkens opnieuw blijft gedetecteerd worden door het model en maakt het overzichtelijker voor het soc-team. 

{\small
\begin{lstlisting}[style=custompython,caption={Deduplicatie mechanisme},label={lst:dlp-selenium}, captionpos=b, basicstyle=\small\ttfamily]
    if not features_df.empty:
    scored_df = score_features(features_df, model, scaler)
    
    now_ts = time.time()
    anomalous_df = scored_df[scored_df["is_anomaly"]].copy()
    
    if not anomalous_df.empty:
        best_per_ip = anomalous_df.loc[
            anomalous_df.groupby("cliIP")["anomaly_score"].idxmin()
        ]
        new_anomalies = best_per_ip[
            ~best_per_ip["cliIP"].apply(
                lambda ip: (now_ts - alerted_ips.get(ip, 0)) < window_size
            )
        ]
    else:
        new_anomalies = pd.DataFrame()
\end{lstlisting}
}

\subsection{Classificatie van anomalieën} 
\label{subsec:classificatie-anomalieen}

Naast de binaire classificatie (anomalie of normaal) voert de pipeline een regelgebaseerde post-classificatie uit die elke gedetecteerde anomalie categoriseert op basis van de onderliggende feature-waarden. De volgende categorieën worden onderscheiden:

\begin{itemize}
    \item \textbf{\textit{HIGH\_RPS}}: meer dan 10 requests per seconde
    \item \textbf{\textit{HIGH\_VOLUME}}: meer dan 200 requests in het venster
    \item \textbf{\textit{HIGH\_4XX}}: meer dan 30\  van de requests resulteren in 4xx-statuscodes
    \item \textbf{\textit{HIGH\_5XX}}: meer dan 10\  van de requests resulteren in 5xx-statuscodes
    \item \textbf{\textit{ENDPOINT\_HAMMERING}}: een path diversity ratio lager dan 0.05 bij meer dan 10 requests
    \item \textbf{\textit{HIGH\_DATA\_TRANSFER}}: meer dan 5~MB aan uitgaande data
    \item \textbf{\textit{AGGRESSIVE\_CRAWL}}: meer dan 100 unieke paden zonder referrer
    \item \textbf{\textit{NO\_REFERER}}: geen enkele referrer bij meer dan 20 requests
    \item \textbf{\textit{UNIFORM\_TIMING}}: een standaardafwijking van de inter-request-tijd kleiner dan 0.01 seconde bij meer dan 10 requests
    \item \textbf{\textit{UA\_INCONSISTENCY}}: meer dan 2 verschillende browsers of besturingssystemen voor één IP
    \item \textbf{\textit{AKAMAI\_MISSED}}: geen enkele Akamai-regel getriggerd (relevant voor verkeer waar Akamai-context beschikbaar is)
\end{itemize}

Een anomalie kan meerdere categorieën tegelijk krijgen, wat het SOC-team helpt om snel te begrijpen welk type verdacht gedrag is gedetecteerd. 

\subsection{classificatie van de ernst van de anomalie} 
\label{subsec:severiteitsbepaling}

Op basis van de anomaliescore wordt een classificatie toegekend hoe ernstig de afwijking is:

\begin{itemize}
    \item \textbf{Critical}: anomaliescore lager dan $-0.08$
    \item \textbf{High}: anomaliescore tussen $-0.08$ en $-0.04$
    \item \textbf{Medium}: anomaliescore tussen $-0.04$ en $-0.01$
    \item \textbf{Low}: anomaliescore hoger dan $-0.01$
\end{itemize}

Hoe lager de anomaliescore, hoe minder splitsingen in de decision trees nodig waren om het datapunt te isoleren en hoe meer het afwijkt van het normale gedragspatroon.

\subsection{Monitoring en visualisatie via Grafana} 
\label{subsec:monitoring-grafana}

De detectieresultaten worden beschikbaar gesteld door het creeëren van een API endpoint die door Grafana opgehaald kan worden. 

\textbf{JSON API.} Een ingebouwde Flask-webserver op poort 8001 biedt een REST API met diverse endpoints die specifiek zijn ontworpen voor integratie met de Grafana Infinity-plugin. Het \textit{/api/anomalies}-endpoint retourneert de volledige lijst van recente anomalieën, \textit{/api/anomalies/summary} groepeert deze per IP-adres met de slechtste score en samengevoegde categorieën, \textit{/api/categories} geeft een verdeling per anomaliecategorie en \textit{/api/severity} biedt een overzicht per severiteitsniveau. Deze API wordt bij het herstarten van de pipeline automatisch gevuld vanuit een persistent CSV-bestand, zodat eerder gedetecteerde anomalieën niet verloren gaan en zichtbaar zijn voor het soc-team in het grafana dashboard.

\subsection{Persistente opslag} 
\label{subsec:persistente-opslag}

Alle gedetecteerde anomalieën worden doorlopend weggeschreven naar een CSV-bestand (\textit{anomalies\_detected.csv}). Bij het opnieuw opstarten van de pipeline wordt dit bestand automatisch ingelezen, zodat de JSON API de volledige anomalie-historie kan presenteren aan Grafana.

\subsection{Configuratie en operationeel beheer} 
\label{subsec:configuratie-operationeel-beheer}

De volledige pipeline wordt aangestuurd via een enkel YAML-configuratiebestand (\textit{config.yaml}) dat de volgende aspecten centraliseert: de Elasticsearch-verbindingsgegevens, polling-parameters (interval en batchgrootte), de venster-instellingen (venstergrootte, verschuivingsstap en minimum requestdrempel), filterregels voor monitoring-verkeers.

De pipeline ondersteunt een graceful shutdown via \textit{SIGINT}- en \textit{SIGTERM}-signalen, wat ervoor zorgt dat de lopende pollcyclus wordt afgerond voordat het proces stopt. Dit is essentieel in een productieomgeving waar het proces wordt beheerd via systemd of een vergelijkbaar init-systeem.
%%=============================================================================
%% Evaluatie
%%=============================================================================

\chapter{\IfLanguageName{dutch}{Evaluatie}{Evaluation}}%
\label{ch:evaluatie}

\section{Evaluatiemethodiek}
\label{sec:evaluatiemethodiek}

De evaluatie van een unsupervised anomaliedetectiemodel verschilt van de evaluatie van een supervised model. Bij supervised learning beschikt de onderzoeker over een volledig gelabelde dataset waarmee klassieke metrics zoals accuracy, precision, recall en F1-score rechtstreeks berekend kunnen worden. Bij het hier geïmplementeerde Isolation Forest model is er geen beschikbare grondwaarde: de Akamai access logs bevatten geen betrouwbare classificatie die voor elk request aangeeft of het om normaal of afwijkend gedrag gaat. Het \textit{aapRule}-veld van Akamai geeft weliswaar aan welke beveiligingsregels zijn getriggerd, maar dit is een indicator van regelgebaseerde detectie doot het akamai systeem en geen grondwaarheid voor gedragsmatige anomalieën.

Om toch een gefundeerde evaluatie te bieden, wordt een driedelige evaluatiemethodiek gebruikt die de prestaties van het model vanuit verschillende invalshoeken analyseert. Het eerste onderdeel is een technische vergelijking van de modelvoorspellingen met de bestaande Akamai-detectie. Het tweede onderdeel is een handmatige validatie van een subset van de gedetecteerde anomalieën door een security-analist. Het derde onderdeel is een operationele beoordeling van de bruikbaarheid van het systeem binnen een SOC-context.

\section{Technische evaluatie}
\label{sec:technische-evaluatie}

\subsection{Vergelijking met Akamai op de trainingsdata}

Om de positionering van het model ten opzichte van de bestaande beveiligingsinfrastructuur te bepalen, worden de voorspellingen van het Isolation Forest model vergeleken met de informatie uit het \textit{aapRule}-veld van Akamai. Per feature-vector wordt bepaald of de meerderheid van de onderliggende requests binnen het window door ten minste één Akamai-regel was gevlagd (\textit{akamai\_has\_rule\_ratio} groter dan 0,5). Deze vergelijking is uitgevoerd op de volledige trainingsdata van 96.289 feature-vectoren en levert vier categorieën op, weergegeven in Tabel~\ref{tab:model-vs-akamai}.

\begin{table}[h]
\centering
\caption{Vergelijking van modelvoorspellingen met Akamai-detectie op de trainingsdata (contamination 0,10).}
\label{tab:model-vs-akamai}
\begin{tabular}{l r r p{6cm}}
\hline
\textbf{Categorie} & \textbf{Aantal} & \textbf{\%} & \textbf{Interpretatie} \\
\hline
Beide vlaggen & 5.598 & 5,8\% & Model bevestigt Akamai's detectie \\
Alleen model & 4.031 & 4,2\% & Model detecteert wat Akamai mist \\
Alleen Akamai & 73.225 & 76,0\% & Regelgebaseerde detectie zonder gedragsafwijking \\
Geen van beide & 13.435 & 14,0\% & Schoon verkeer \\
\hline
\end{tabular}
\end{table}

De resultaten tonen dat het model 5.598 feature-vectoren als anomalie classificeert die ook door Akamai zijn gevlagd, wat bevestigt dat het model in staat is om verkeer te herkennen dat de bestaande rule-engine als verdacht beschouwt. De 4.031 vectoren die uitsluitend door het model zijn gedetecteerd vormen de primaire meerwaarde van het proof of concept: dit is verkeer dat door de Akamai rule-engine glipt maar dat het ML-model wél als gedragsmatig afwijkend classificeert.

De 73.225 vectoren die alleen door Akamai zijn gevlagd maar niet door het model vormen geen gemiste detecties. Akamai vlagt verkeer op basis van bekende signatures en rule-IDs, zoals \textit{BOT-BROWSER-IMPERSONATOR}. Dit verkeer voldoet aan de voorwaarden van een specifieke regel, maar vertoont qua gedrag geen statistisch afwijkend patroon ten opzichte van het normale verkeer dat het model heeft geleerd. Dit benadrukt dat het ML-model en de rule-engine complementaire detectiemethoden zijn.

\subsection{Detectie op verkeer zonder Akamai-regels}

Om specifiek de meerwaarde van het model als aanvulling op Akamai te evalueren, is het getrainde model toegepast op een dataset van 85.501 logrecords waarvoor het \textit{aapRule}-veld leeg was. Dit zijn requests die de volledige Akamai rule-engine hebben doorlopen zonder dat er een WAF-regel, bot-regel of custom regel is getriggerd. Na preprocessing en feature engineering resulteerde dit in 22.121 feature-vectoren van 3.001 unieke IP-adressen. Het model detecteerde 2.083 anomalieën (9,4\%), weergegeven in Tabel~\ref{tab:aap-empty}.

\begin{table}[h]
\centering
\caption{Resultaten van het model op verkeer zonder Akamai-detectie.}
\label{tab:aap-empty}
\begin{tabular}{l r}
\hline
\textbf{Metriek} & \textbf{Waarde} \\
\hline
Totaal gescorde windows & 22.121 \\
Unieke IP-adressen & 3.001 \\
Gedetecteerde anomalieën & 2.083 (9,4\%) \\
Normaal verkeer & 20.038 (90,6\%) \\
\hline
\end{tabular}
\end{table}

Aangezien deze logs per definitie geen Akamai-regel hebben getriggerd, zijn alle detecties verkeer dat de bestaande beveiligingsinfrastructuur heeft doorlaten. De meest voorkomende anomaliecategorieën zijn \textit{HIGH\_5XX} (67,0\%), \textit{UA\_INCONSISTENCY} (23,2\%), \textit{NO\_REFERER} (12,3\%) en \textit{ENDPOINT\_HAMMERING} (6,8\%). De dominantie van \textit{HIGH\_5XX} is opmerkelijk: dit zijn requests die server errors veroorzaken, wat kan wijzen op fuzz testing of applicatie-misbruik. Deze requests worden niet door Akamai gevlagd omdat ze geen bekende aanvalssignatuur bevatten, maar het gedragspatroon wordt door het model wél als afwijkend herkend.

\subsection{Vergelijking met gelijkaardig onderzoek}

De trainingsresultaten worden vergeleken met het onderzoek van Chua et al.~\autocite{Chua2024}, die eveneens Isolation Forest implementeerden voor anomaliedetectie in webverkeer. Tabel~\ref{tab:vergelijking-chua} toont de vergelijking van de modelparameters en de resultaten.

\begin{table}[h]
\centering
\caption{Vergelijking van het proof of concept met Chua et al. (2024).}
\label{tab:vergelijking-chua}
\begin{tabular}{l l l}
\hline
\textbf{Aspect} & \textbf{Dit onderzoek} & \textbf{Chua et al. (2024)} \\
\hline
Dataset & Akamai access logs & Nginx server logs \\
Omvang trainingsdata & 96.289 feature-vectoren & 2.591.287 records \\
Contamination & 0,10 & 0,03 \\
Features & 26 gedragsfeatures & 13 features (incl. IoC) \\
Feature-aanpak & Sliding window per IP & Per request \\
Precision (breed) & 93,7\% & 95\% \\
False positive rate & 6,3\% & 4,5\% \\
Recall & Niet berekend & 90\% \\
Near-real-time & Ja (30s polling) & Nee (batch) \\
\hline
\end{tabular}
\end{table}

De precision van dit onderzoek (93,7\%) is vergelijkbaar met die van Chua et al. (95\%), ondanks de verschillende contexten. Het verschil in false positive rate (6,3\% versus 4,5\%) kan verklaard worden door de hogere contamination parameter (0,10 versus 0,03) en het feit dat dit onderzoek de evaluatie op IP-niveau uitvoert in plaats van op request-niveau. Een belangrijk verschil is dat dit onderzoek een near-real-time pipeline implementeert met een effectieve detectielatentie van circa 30 seconden, terwijl Chua et al. uitsluitend batch-verwerking toepassen.

\section{Handmatige validatie}
\label{sec:handmatige-validatie}

\subsection{Methodiek}

Om de precision van het model te bepalen en inzicht te krijgen in de aard van de gedetecteerde anomalieën, is een handmatige validatie uitgevoerd op de unieke IP-adressen die door het model als anomaal zijn geclassificeerd gedurende de near-real-time detectieperiode. Voor elk IP zijn de volledige requestlogs opgehaald uit Elasticsearch en is een gedetailleerd gedragsrapport gegenereerd. Dit rapport bevat de volledige requestgeschiedenis, de statuscodeverdeling, de gebruikte HTTP-methodes, de meest bezochte paden, de user-agent informatie, een timinganalyse, het datavolume, en de Akamai-detectiestatus.

Om een representatief beeld te verkrijgen van de modelprestaties is een gestratificeerde steekproef genomen over de vier severity-niveaus (critical, high, medium, low). Dit voorkomt dat de evaluatie vertekend wordt door uitsluitend de meest extreme anomalieën te beoordelen, en geeft inzicht in de betrouwbaarheid van het model bij grensgevallen.

Op basis van de gedragsrapporten is elk IP geclassificeerd in een van vier categorieën:

\begin{itemize}
    \item \textbf{True Positive --- Malicious:} verkeer dat aantoonbaar kwaadaardig is, zoals vulnerability scanners, brute force pogingen en agressieve scrapers.
    \item \textbf{True Positive --- Suspicious Bot:} geautomatiseerd verkeer dat verdacht is maar waarvan de intentie niet eenduidig als kwaadaardig kan worden vastgesteld.
    \item \textbf{True Positive --- Benign Bot:} goedaardig geautomatiseerd verkeer, zoals bekende SEO-crawlers, e-mail clients en monitoring-tools.
    \item \textbf{False Positive:} IP-adressen die door het model als anomaal zijn geclassificeerd maar bij nader onderzoek normaal gebruikersgedrag vertonen.
\end{itemize}

De classificatie als true positive of false positive is niet uitsluitend gebaseerd op de vraag of het verkeer kwaadaardig is. Het model is ontworpen om gedragsafwijkingen te detecteren, niet om onderscheid te maken tussen goedaardige en kwaadaardige bots. Een goedaardige bot die daadwerkelijk afwijkend gedrag vertoont is een terechte detectie (true positive), omdat het model correct heeft geïdentificeerd dat het gedrag afwijkt van het normale patroon. Alleen verkeer dat geen daadwerkelijke gedragsafwijking vertoont wordt als false positive beschouwd.

\subsection{Resultaten van de handmatige validatie}

% ============================================================
% PLACEHOLDER: Vul hier de resultaten in van de validatie door de Evolane-analist
% ============================================================

De handmatige validatie is uitgevoerd op \textit{[AANTAL]} unieke IP-adressen, verdeeld over de vier severity-niveaus. De resultaten zijn weergegeven in Tabel~\ref{tab:handmatige-validatie}.

\begin{table}[h]
\centering
\caption{Classificatie van de geëvalueerde IP-adressen door de security-analist.}
\label{tab:handmatige-validatie}
\begin{tabular}{l r r p{5cm}}
\hline
\textbf{Classificatie} & \textbf{Aantal} & \textbf{\%} & \textbf{Beschrijving} \\
\hline
True Positive --- Malicious & \textit{[X]} & \textit{[X\%]} & Scanners, brute force, webshell zoeken \\
True Positive --- Suspicious Bot & \textit{[X]} & \textit{[X\%]} & Verdachte bots, UA-rotatie \\
True Positive --- Benign Bot & \textit{[X]} & \textit{[X\%]} & SEO-crawlers, e-mail clients \\
False Positive & \textit{[X]} & \textit{[X\%]} & Normaal verkeer onterecht gevlagd \\
\hline
\end{tabular}
\end{table}

De handmatige validatie levert twee precisiemetrieken op. Wanneer elke terechte gedragsafwijking als true positive wordt beschouwd, bedraagt de precision \textit{[X\%]}. Wanneer uitsluitend kwaadaardig verkeer als true positive wordt geteld, bedraagt de precision \textit{[X\%]}. De false positive rate bedraagt \textit{[X\%]}.

% ============================================================
% PLACEHOLDER: Voeg hier eventueel 2-3 concrete voorbeelden toe
% van opvallende detecties die de analist heeft beoordeeld
% ============================================================

\section{Operationele beoordeling}
\label{sec:operationele-beoordeling}

Naast de technische evaluatie is ook de operationele bruikbaarheid van het systeem beoordeeld. Deze beoordeling richt zich op de vraag of het proof of concept integreerbaar is met de bestaande securityomgeving van Evolane en of het in de praktijk bruikbare informatie oplevert voor het SOC-team.

\subsection{Near-real-time detectieperformantie}

De pipeline is geïmplementeerd als een continu draaiend proces op een Ubuntu VM in de Google Cloud omgeving. De volledige verwerkingscyclus van het ophalen van logs uit Elasticsearch tot het scoren met het Isolation Forest model duurt doorgaans minder dan twee seconden. Met een polling-interval van 30 seconden resulteert dit in een effectieve detectielatentie van circa 30 seconden. Dit betekent dat afwijkend verkeer binnen een halve minuut na het plaatsvinden kan worden gedetecteerd en gerapporteerd.

De pipeline is ontworpen voor langdurige, ononderbroken uitvoering. Een graceful shutdown mechanisme zorgt ervoor dat het proces bij een stopverzoek de huidige cyclus afrondt. Een rolling buffer van tien minuten aan preprocessed logs garandeert dat er bij fluctuerend verkeersvolume altijd voldoende data beschikbaar is voor volledige sliding windows. Een deduplicatiemechanisme op IP-niveau met een cooldown-periode voorkomt herhaalde alerts voor hetzelfde IP binnen een kort tijdsbestek.

\subsection{Persistentie en beschikbaarheid}

Gedetecteerde anomalieën worden persistent opgeslagen in een CSV-bestand. Bij een herstart van de pipeline of de VM worden eerder gedetecteerde anomalieën automatisch hersteld vanuit dit bestand, zodat het Grafana-dashboard direct de historische detecties toont. De pipeline is geconfigureerd als een systemd-service die automatisch opstart bij een reboot van de VM en automatisch herstart bij een onverwachte crash.

\subsection{Visualisatie en integratie}

De detectieresultaten worden ontsloten via een JSON API die door de Grafana Infinity plugin wordt bevraagd. Het Grafana-dashboard biedt een overzicht van de gedetecteerde anomalieën met de mogelijkheid om door te klikken naar de onderliggende Elasticsearch-logs per IP-adres. Dit stelt een security-analist in staat om vanuit het anomalie-overzicht direct de ruwe logs te inspecteren en een beoordeling te maken van de gedetecteerde afwijking.

\subsection{Beperkingen}

Het systeem kent een aantal operationele beperkingen. Ten eerste is de vergelijking met Akamai-detectie op window-niveau uitgevoerd in plaats van op IP-niveau, waardoor de \textit{akamai\_detected} indicator niet altijd het volledige beeld geeft van wat Akamai over een specifiek IP weet. Ten tweede is het model getraind op een dataset die reeds veel botverkeer bevatte, waardoor het model bepaalde botpatronen als \textit{normaal} heeft geleerd. Ten derde kan recall niet berekend worden zonder een volledig gelabelde referentiedataset, wat een beperking is van de evaluatiemethode.



%%=============================================================================
%% Resultaten
%%=============================================================================

\chapter{\IfLanguageName{dutch}{Resultaten}{Results}}%
\label{ch:resultaten}

\section{Inleiding}
\label{sec:inleiding-resultaten}

Dit hoofdstuk presenteert de concrete resultaten van het anomaliedetectiesysteem. Eerst worden de trainingsresultaten besproken, gevolgd door de resultaten van de near-real-time detectie en de uitkomsten van de handmatige validatie. Tot slot wordt een samenvatting gegeven van alle evaluatiemetrieken.

\section{Trainingsresultaten}
\label{sec:trainingsresultaten}

Het Isolation Forest model is getraind op een historische dataset van 626.263 Akamai access logs, verzameld over een periode van ongeveer vijftien dagen. Na preprocessing, waarbij 53.443 monitoringrequests en 40 health checks zijn verwijderd, resteerden 572.780 bruikbare records. De feature engineering over sliding windows van vijf minuten resulteerde in 96.289 feature-vectoren van 3.556 unieke IP-adressen. De trainingsresultaten zijn samengevat in Tabel~\ref{tab:trainingsresultaten}.

\begin{table}[h]
\centering
\caption{Samenvatting van de trainingsresultaten.}
\label{tab:trainingsresultaten}
\begin{tabular}{l r}
\hline
\textbf{Metriek} & \textbf{Waarde} \\
\hline
Ruwe logrecords & 626.263 \\
Na preprocessing & 572.780 \\
Feature-vectoren & 96.289 \\
Unieke IP-adressen & 3.556 \\
Contamination parameter & 0,10 \\
Aantal decision trees & 200 \\
Trainingstijd & 1,5 seconden \\
Gedetecteerde anomalieën & 9.629 (10,0\%) \\
Normaal verkeer & 86.660 (90,0\%) \\
Anomaly score bereik & [-0,1785 ; 0,1592] \\
\hline
\end{tabular}
\end{table}

De keuze voor een contamination van 0,10 is het resultaat van een iteratief proces. Een eerste training met een contamination van 0,05 resulteerde in 4.813 anomalieën waarvan 1.899 nieuwe vondsten ten opzichte van Akamai. Na verhoging naar 0,10 steeg het aantal nieuwe vondsten naar 4.031, terwijl de overlap met Akamai eveneens toenam van 2.914 naar 5.598. Deze dubbele stijging bevestigt dat de extra detecties niet uitsluitend ruis zijn, maar dat het model bij een hogere sensitiviteit ook meer bevestigt van wat Akamai al detecteert.

\section{Resultaten vergelijking model en Akamai}
\label{sec:resultaten-model-akamai}

De vergelijking tussen de modelvoorspellingen en de Akamai-detectie op de trainingsdata toont de complementariteit van beide systemen. Het model vindt 4.031 anomalieën die Akamai mist, terwijl Akamai 73.225 gevallen vlagt die het model als gedragsmatig normaal beschouwt. De 5.598 gevallen waarin beide systemen overeenstemmen bevestigen dat het model in staat is relevante anomalieën te herkennen.

Bij de toepassing op de 85.501 logs zonder Akamai-regels detecteert het model 2.083 aanvullende anomalieën. Dit zijn IP-adressen die de volledige Akamai rule-engine hebben doorlopen zonder enige detectie, maar die op basis van hun gedragsprofiel wel afwijken van het normale patroon. De voornaamste categorieën zijn verkeer dat server errors veroorzaakt (67\%), user-agent inconsistenties (23\%) en verkeer zonder referer header (12\%).

\section{Resultaten near-real-time detectie}
\label{sec:resultaten-near-real-time}

Gedurende de near-real-time detectieperiode heeft de pipeline \textit{[AANTAL]} unieke IP-adressen als anomaal geclassificeerd. De verdeling over de severity-niveaus is weergegeven in Tabel~\ref{tab:severity-verdeling}.

\begin{table}[h]
\centering
\caption{Verdeling van gedetecteerde anomalieën per severity-niveau.}
\label{tab:severity-verdeling}
\begin{tabular}{l r r}
\hline
\textbf{Severity} & \textbf{Aantal IPs} & \textbf{Percentage} \\
\hline
Critical (score $<$ -0,08) & \textit{[X]} & \textit{[X\%]} \\
High (score -0,08 tot -0,04) & \textit{[X]} & \textit{[X\%]} \\
Medium (score -0,04 tot -0,01) & \textit{[X]} & \textit{[X\%]} \\
Low (score -0,01 tot 0,00) & \textit{[X]} & \textit{[X\%]} \\
\hline
\end{tabular}
\end{table}

% ============================================================
% PLACEHOLDER: Vul de severity aantallen in vanuit je Grafana dashboard
% ============================================================

\subsection{Voorbeelden van gedetecteerde anomalieën}

Om de werking van het model te illustreren worden drie representatieve detecties besproken die de verschillende typen anomalieën tonen die het systeem identificeert.

% ============================================================
% PLACEHOLDER: Kies 3 concrete IPs uit je investigation reports
% en beschrijf ze hieronder. Suggestie:
% 1. Een duidelijk kwaadaardige (bv. de WordPress scanner)
% 2. Een verdachte bot (bv. de rb_dzc35632 bezoekers)
% 3. Een goedaardige bot (bv. Google Gmail Image Proxy)
% ============================================================

\textbf{Voorbeeld 1: Vulnerability scanner.} Het IP-adres \textit{[IP]} werd gedetecteerd met een anomaly score van \textit{[SCORE]} (severity: critical). In een tijdsvenster van vijf minuten werden \textit{[AANTAL]} requests uitgevoerd naar \textit{[AANTAL]} unieke paden, waaronder verdachte bestanden zoals \textit{[PADEN]}. De error rate bedroeg \textit{[X\%]} en er was geen user-agent of referer header aanwezig. Het model classificeerde dit als \textit{[CATEGORIEËN]}. Dit IP was \textit{[wel/niet]} door Akamai gedetecteerd.

\textbf{Voorbeeld 2: Verdachte bot.} Het IP-adres \textit{[IP]} werd gedetecteerd met een anomaly score van \textit{[SCORE]}. \textit{[Beschrijf het gedrag en waarom het verdacht is.]}

\textbf{Voorbeeld 3: Goedaardige bot.} Het IP-adres \textit{[IP]} werd gedetecteerd met een anomaly score van \textit{[SCORE]}. \textit{[Beschrijf het gedrag en leg uit waarom het een true positive maar goedaardig is.]}

\section{Resultaten handmatige validatie}
\label{sec:resultaten-handmatige-validatie}

% ============================================================
% PLACEHOLDER: Vul in na de validatie door de Evolane-analist
% ============================================================

De gestratificeerde steekproef van \textit{[AANTAL]} IP-adressen is voorgelegd aan een security-analist van Evolane ter handmatige classificatie. De resultaten zijn weergegeven in Tabel~\ref{tab:resultaten-validatie}.

\begin{table}[h]
\centering
\caption{Resultaten van de handmatige validatie.}
\label{tab:resultaten-validatie}
\begin{tabular}{l r r}
\hline
\textbf{Classificatie} & \textbf{Aantal} & \textbf{Percentage} \\
\hline
True Positive --- Malicious & \textit{[X]} & \textit{[X\%]} \\
True Positive --- Suspicious Bot & \textit{[X]} & \textit{[X\%]} \\
True Positive --- Benign Bot & \textit{[X]} & \textit{[X\%]} \\
False Positive & \textit{[X]} & \textit{[X\%]} \\
\hline
\textbf{Totaal} & \textit{[X]} & 100\% \\
\hline
\end{tabular}
\end{table}

Op basis van deze validatie zijn de volgende precisiemetrieken berekend:

\begin{itemize}
    \item \textbf{Precision (breed):} het percentage detecties dat een daadwerkelijke gedragsafwijking betreft (alle categorieën behalve false positive): \textit{[X\%]}.
    \item \textbf{Precision (strikt):} het percentage detecties dat aantoonbaar kwaadaardig verkeer betreft (alleen malicious): \textit{[X\%]}.
    \item \textbf{False positive rate:} het percentage detecties dat normaal gebruikersgedrag betreft: \textit{[X\%]}.
\end{itemize}

\section{Samenvatting van de resultaten}
\label{sec:samenvatting-resultaten}

Tabel~\ref{tab:samenvatting-alle-resultaten} geeft een overzicht van alle evaluatieresultaten.

\begin{table}[h]
\centering
\caption{Samenvatting van alle evaluatieresultaten.}
\label{tab:samenvatting-alle-resultaten}
\begin{tabular}{l l r}
\hline
\textbf{Evaluatieonderdeel} & \textbf{Metriek} & \textbf{Waarde} \\
\hline
Model vs Akamai & Overlap (beide vlaggen) & 5.598 (5,8\%) \\
Model vs Akamai & Nieuwe detecties (model only) & 4.031 (4,2\%) \\
Ongedetecteerd verkeer & Anomalieën zonder aapRule & 2.083 / 22.121 (9,4\%) \\
Handmatige validatie & Precision (breed) & \textit{[X\%]} \\
Handmatige validatie & Precision (strikt) & \textit{[X\%]} \\
Handmatige validatie & False positive rate & \textit{[X\%]} \\
Handmatige validatie & Geëvalueerde IPs & \textit{[X]} \\
Near-real-time & Detectielatentie & $\sim$30 seconden \\
Near-real-time & Verwerkingstijd per cyclus & $<$ 2 seconden \\
\hline
\end{tabular}
\end{table}

\subsection{Beperking: recall}

Een belangrijke beperking is dat recall niet berekend kan worden. Recall vereist kennis van het totale aantal werkelijke anomalieën in de dataset, inclusief de anomalieën die het model niet heeft gedetecteerd. Zonder een volledig gelabelde dataset is het onmogelijk om te bepalen hoeveel afwijkend verkeer het model heeft gemist. Vergelijkbaar onderzoek van Chua et al.~\autocite{Chua2024} bereikte een recall van 90\% met Isolation Forest op webverkeer, wat een indicatie biedt maar niet rechtstreeks overdraagbaar is naar de context van dit onderzoek. Het labelen van een representatieve subset van de Akamai logs door het security team van Evolane zou de berekening van recall mogelijk maken, maar valt buiten de scope van dit proof of concept.


% Voeg hier je eigen hoofdstukken toe die de ``corpus'' van je bachelorproef
% vormen. De structuur en titels hangen af van je eigen onderzoek. Je kan bv.
% elke fase in je onderzoek in een apart hoofdstuk bespreken.

%\input{...}
%\input{...}
%...

\input{conclusie}

%---------- Bijlagen -----------------------------------------------------------

\appendix

\chapter{Onderzoeksvoorstel}

Het onderwerp van deze bachelorproef is gebaseerd op een onderzoeksvoorstel dat vooraf werd beoordeeld door de promotor. Dat voorstel is opgenomen in deze bijlage.

%% TODO: 
%\section*{Samenvatting}

% Kopieer en plak hier de samenvatting (abstract) van je onderzoeksvoorstel.

% Verwijzing naar het bestand met de inhoud van het onderzoeksvoorstel
\input{../voorstel/voorstel-inhoud}

%%---------- Andere bijlagen --------------------------------------------------
% TODO: Voeg hier eventuele andere bijlagen toe. Bv. als je deze BP voor de
% tweede keer indient, een overzicht van de verbeteringen t.o.v. het origineel.
%\input{...}

%%---------- Backmatter, referentielijst ---------------------------------------

\backmatter{}

\setlength\bibitemsep{2pt} %% Add Some space between the bibliograpy entries
\printbibliography[heading=bibintoc]

\end{document}